% Options for packages loaded elsewhere
\PassOptionsToPackage{unicode}{hyperref}
\PassOptionsToPackage{hyphens}{url}
\PassOptionsToPackage{dvipsnames,svgnames,x11names}{xcolor}
%
\documentclass[
  a4paper,
  13pt]{extreport}
\usepackage{amsmath,amssymb}
\usepackage{setspace}
\usepackage{iftex}
\ifPDFTeX
  \usepackage[T1]{fontenc}
  \usepackage[utf8]{inputenc}
  \usepackage{textcomp} % provide euro and other symbols
\else % if luatex or xetex
  \usepackage{unicode-math} % this also loads fontspec
  \defaultfontfeatures{Scale=MatchLowercase}
  \defaultfontfeatures[\rmfamily]{Ligatures=TeX,Scale=1}
\fi
\usepackage{lmodern}
\ifPDFTeX\else
  % xetex/luatex font selection
  \setmainfont[]{TeX Gyre Termes}
\fi
% Use upquote if available, for straight quotes in verbatim environments
\IfFileExists{upquote.sty}{\usepackage{upquote}}{}
\IfFileExists{microtype.sty}{% use microtype if available
  \usepackage[]{microtype}
  \UseMicrotypeSet[protrusion]{basicmath} % disable protrusion for tt fonts
}{}
\usepackage{xcolor}
\usepackage[left=3cm,right=2cm,top=2cm,bottom=2cm,headsep=1cm,footskip=1cm]{geometry}
\usepackage{longtable,booktabs,array}
\usepackage{calc} % for calculating minipage widths
% Correct order of tables after \paragraph or \subparagraph
\usepackage{etoolbox}
\makeatletter
\patchcmd\longtable{\par}{\if@noskipsec\mbox{}\fi\par}{}{}
\makeatother
% Allow footnotes in longtable head/foot
\IfFileExists{footnotehyper.sty}{\usepackage{footnotehyper}}{\usepackage{footnote}}
\makesavenoteenv{longtable}
\usepackage{graphicx}
\makeatletter
\def\maxwidth{\ifdim\Gin@nat@width>\linewidth\linewidth\else\Gin@nat@width\fi}
\def\maxheight{\ifdim\Gin@nat@height>\textheight\textheight\else\Gin@nat@height\fi}
\makeatother
% Scale images if necessary, so that they will not overflow the page
% margins by default, and it is still possible to overwrite the defaults
% using explicit options in \includegraphics[width, height, ...]{}
\setkeys{Gin}{width=\maxwidth,height=\maxheight,keepaspectratio}
% Set default figure placement to htbp
\makeatletter
\def\fps@figure{htbp}
\makeatother
\setlength{\emergencystretch}{3em} % prevent overfull lines
\providecommand{\tightlist}{%
  \setlength{\itemsep}{0pt}\setlength{\parskip}{0pt}}
\setcounter{secnumdepth}{-\maxdimen} % remove section numbering
\ifLuaTeX
\usepackage[bidi=basic]{babel}
\else
\usepackage[bidi=default]{babel}
\fi
\babelprovide[main,import]{vietnamese}
\ifPDFTeX
\else
\babelfont{rm}[]{TeX Gyre Termes}
\fi
% get rid of language-specific shorthands (see #6817):
\let\LanguageShortHands\languageshorthands
\def\languageshorthands#1{}
% ctu_xelatex_header.tex — header-includes cho build LaTeX CTU (xelatex).
% Ánh xạ các ký tự Unicode mà TeX Gyre Termes thiếu glyph (chỉ số dưới/trên,
% ký tự kẻ hộp) sang lệnh LaTeX tương đương để không bị rớt glyph khi biên dịch.
\usepackage{newunicodechar}
\usepackage{pifont}
\usepackage{amssymb}
\newunicodechar{✓}{\ding{51}}
\newunicodechar{✗}{\ding{55}}
\newunicodechar{≥}{\ensuremath{\geq}}
\newunicodechar{≤}{\ensuremath{\leq}}
\newunicodechar{≈}{\ensuremath{\approx}}
\newunicodechar{↔}{\ensuremath{\leftrightarrow}}
\newunicodechar{₀}{\textsubscript{0}}
\newunicodechar{₁}{\textsubscript{1}}
\newunicodechar{₂}{\textsubscript{2}}
\newunicodechar{₃}{\textsubscript{3}}
\newunicodechar{₄}{\textsubscript{4}}
\newunicodechar{₅}{\textsubscript{5}}
\newunicodechar{₆}{\textsubscript{6}}
\newunicodechar{₇}{\textsubscript{7}}
\newunicodechar{₈}{\textsubscript{8}}
\newunicodechar{₉}{\textsubscript{9}}
\newunicodechar{²}{\textsuperscript{2}}
\newunicodechar{³}{\textsuperscript{3}}
\newunicodechar{─}{\rule[0.5ex]{0.9em}{0.4pt}}

% --- Bảng rộng: thu nhỏ cỡ chữ + khoảng cột, cho phép ngắt token dài để
%     bảng không tràn lề phải (overfull hbox). Áp dụng cho mọi longtable
%     (pandoc sinh toàn bộ bảng dạng longtable).
\usepackage{etoolbox}
\usepackage{seqsplit}      % ngắt được token rất dài (mã biến, số dài)
% Lưu ý: KHÔNG đổi \tabcolsep — pandoc tính bề rộng cột theo công thức
% (\columnwidth - 4\tabcolsep)*frac (tham chiếu \tabcolsep trực tiếp), nên
% giảm \tabcolsep lại làm cột RỘNG hơn và tràn nặng hơn.
\AtBeginEnvironment{longtable}{%
  \footnotesize
  \emergencystretch=3em
  \hbadness=10000          % không cảnh báo underfull do RaggedRight
}
% Cho phép TeX bẻ dòng linh hoạt hơn trong ô p{} hẹp
\setlength{\emergencystretch}{2em}

% URL/DOI dài: cho phép ngắt ở bất kỳ ký tự nào (tránh tràn lề ở danh mục TLTK).
% Kết hợp với pandoc +autolink_bare_uris để URL trần được bọc \url{}.
\usepackage{xurl}
\ifLuaTeX
  \usepackage{selnolig}  % disable illegal ligatures
\fi
\IfFileExists{bookmark.sty}{\usepackage{bookmark}}{\usepackage{hyperref}}
\IfFileExists{xurl.sty}{\usepackage{xurl}}{} % add URL line breaks if available
\urlstyle{same}
\hypersetup{
  pdflang={vi},
  colorlinks=true,
  linkcolor={black},
  filecolor={Maroon},
  citecolor={black},
  urlcolor={blue},
  pdfcreator={LaTeX via pandoc}}

\author{}
\date{}

\begin{document}

\setstretch{1.2}
\hypertarget{chux1b0ux1a1ng-3-phux1b0ux1a1ng-phuxe1p-nghiuxean-cux1ee9u}{%
\chapter{CHƯƠNG 3: PHƯƠNG PHÁP NGHIÊN
CỨU}\label{chux1b0ux1a1ng-3-phux1b0ux1a1ng-phuxe1p-nghiuxean-cux1ee9u}}

\begin{center}\rule{0.5\linewidth}{0.5pt}\end{center}

\hypertarget{thiux1ebft-kux1ebf-nghiuxean-cux1ee9u-tux1ed5ng-thux1ec3}{%
\section{3.1 Thiết kế nghiên cứu tổng
thể}\label{thiux1ebft-kux1ebf-nghiuxean-cux1ee9u-tux1ed5ng-thux1ec3}}

\hypertarget{loux1ea1i-thiux1ebft-kux1ebf}{%
\subsection{3.1.1 Loại thiết kế}\label{loux1ea1i-thiux1ebft-kux1ebf}}

Luận án sử dụng \textbf{thiết kế hỗn hợp tổng hợp--thực nghiệm} gồm hai
giai đoạn bổ sung cho nhau. Giai đoạn thứ nhất là phân tích tổng hợp tài
liệu để trả lời RQ1, còn giai đoạn thứ hai là phân tích thực nghiệm
doanh nghiệp đa quốc gia để trả lời RQ2--RQ4. Hai giai đoạn đối thoại
hai chiều: phân tích tổng hợp cung cấp cơ sở và gợi ý biến điều tiết cần
kiểm định, còn thực nghiệm cung cấp bằng chứng bổ sung cho phân tích
tổng hợp về bối cảnh châu Á và Pacific.

\textbf{Cơ sở kế thừa}: Thiết kế hỗn hợp kết hợp phân tích tổng hợp và
dữ liệu sơ cấp đã được áp dụng trong quản trị chiến lược và IB
(Borenstein et al., 2009; Hunter \& Schmidt, 2004; Combs et al., 2011).

\textbf{Đóng góp mới}: Áp dụng thiết kế hỗn hợp cho một luận án duy nhất
về quốc tế hóa--hiệu quả hoạt động kinh doanh của doanh nghiệp trong bối
cảnh châu Á và Thái Bình Dương với \textbf{50 nền kinh tế} là thiết kế
hiếm gặp; đa số các luận án IB chỉ làm một trong hai giai đoạn. Quy mô
50 nền kinh tế/88.869 doanh nghiệp là phạm vi địa lý lớn nhất từng có
cho khu vực, mở rộng từ mẫu gộp 17 nước trong nghiên cứu thành phần P1.

\hypertarget{quy-truxecnh-nghiuxean-cux1ee9u}{%
\subsection{3.1.2 Quy trình nghiên
cứu}\label{quy-truxecnh-nghiuxean-cux1ee9u}}

Quy trình đi theo sáu bước: (i) xác định khoảng trống nghiên cứu từ tổng
quan và phân tích tổng hợp trước đây; (ii) xây dựng khung lý thuyết tích
hợp đa tầng và hệ giả thuyết H1--H6; (iii) tiến hành phân tích tổng hợp
cập nhật 1977--2026; (iv) thu thập và chuẩn hóa dữ liệu WBES cho
\textbf{50 nền kinh tế châu Á và Thái Bình Dương} sau khi hòa hợp 3 thế
hệ lược đồ (PICS3, Standardized, BREADY/BEE); (v) ước lượng mô hình thực
nghiệm theo bối cảnh quốc gia và đa quốc gia; (vi) kiểm định độ vững và
tích hợp kết quả để đóng góp lý thuyết. Toàn bộ quy trình và quan hệ bổ
sung hai chiều giữa hai giai đoạn được trình bày ở Hình 3.1.

\includegraphics{figures/fig_3_1_research_process.png}

\textbf{Hình 3.1.} \emph{Quy trình nghiên cứu của luận án.} Mũi tên liền
thể hiện trình tự thực hiện; mũi tên hai chiều thể hiện quan hệ bổ sung
giữa hai giai đoạn (phân tích tổng hợp gợi ý biến điều tiết cần kiểm
định; phân tích thực nghiệm cung cấp bằng chứng châu Á và Thái Bình
Dương cho phân tích tổng hợp). \emph{Nguồn: tác giả xây dựng.}

\hypertarget{giai-ux111oux1ea1n-1-phuxe2n-tuxedch-tux1ed5ng-hux1ee3p-19772026}{%
\section{3.2 Giai đoạn 1: Phân tích tổng hợp
1977--2026}\label{giai-ux111oux1ea1n-1-phuxe2n-tuxedch-tux1ed5ng-hux1ee3p-19772026}}

\hypertarget{cux1a1-sux1edf-kux1ebf-thux1eeba}{%
\subsection{3.2.1 Cơ sở kế
thừa}\label{cux1a1-sux1edf-kux1ebf-thux1eeba}}

phân tích tổng hợp trong IB đã được thiết lập vững chắc qua nhiều thập
kỷ (Hunter \& Schmidt, 2004; Borenstein et al., 2009). Phương pháp chuẩn
bao gồm: xây dựng tiêu chí lựa chọn theo PRISMA (Page et al., 2021),
chuyển dạng hiệu ứng bằng Fisher's \(z\), tính hiệu ứng gộp size theo
hiệu ứng ngẫu nhiên, kiểm định dị biệt bằng \(Q\)-test và \(I^2\), và
đánh giá lệch lạc công bố bằng Egger và Begg-Mazumdar (Egger et al.,
1997; Begg \& Mazumdar, 1994). Các các phân tích tổng hợp trước đây về
quốc tế hóa và hiệu quả (Bausch \& Krist, 2007; Kirca et al., 2012;
Marano et al., 2016; Wu et al., 2022; Arte \& Larimo, 2022) cung cấp quy
trình mẫu mà luận án kế thừa và cập nhật.

\hypertarget{ux111uxf3ng-guxf3p-mux1edbiux111iux1ec1u-chux1ec9nh}{%
\subsection{3.2.2 Đóng góp mới/điều
chỉnh}\label{ux111uxf3ng-guxf3p-mux1edbiux111iux1ec1u-chux1ec9nh}}

\begin{itemize}
\tightlist
\item
  \textbf{Mở rộng coverage time}: từ 1977--2022 (bài ICBEF 2024) sang
  \textbf{1977--2026}, bổ sung các nghiên cứu 2023--2026 liên quan đến
  chuyển đổi số và doanh nghiệp đa quốc gia thị trường mới nổi.
\item
  \textbf{Mở rộng biến điều tiết coverage}: thêm cụm biến điều tiết về
  năng lực số và 6 chế độ con ICRV, chưa được xem xét đầy đủ trong các
  các phân tích tổng hợp trước.
\item
  \textbf{Tập tổng hợp}: K=288 cỡ ảnh hưởng / k=238 nghiên cứu, được xây
  dựng theo chiến lược tìm kiếm hệ thống nhưng có giới hạn (neo theo
  trích dẫn: quét trích dẫn xuôi và ngược từ 5 phân tích tổng hợp mỏ
  neo, bổ sung bằng cơ sở dữ liệu), báo cáo theo biến thể PRISMA 2020
  ``studies identified via other methods''; mở rộng so với bản nền ICBEF
  2024 (K=200 / k=113).
\item
  \textbf{Subgroup theo châu Á và Pacific}: tách riêng nhóm nghiên cứu
  về doanh nghiệp châu Á và Pacific để cung cấp cơ sở trực tiếp cho giai
  đoạn thực nghiệm, đối chiếu với mẫu gộp 17 nước trong nghiên cứu thành
  phần P1.
\end{itemize}

\hypertarget{quy-truxecnh-prisma}{%
\subsection{3.2.3 Quy trình PRISMA}\label{quy-truxecnh-prisma}}

Bốn bước Nhận diện, Sàng lọc, Đủ điều kiện và Đưa vào theo hướng dẫn
PRISMA 2020 (Page et al., 2021). tiêu chí lựa chọn: nghiên cứu thực
nghiệm ở cấp doanh nghiệp, có đo lường quốc tế hóa và hiệu quả hoạt động
kinh doanh của doanh nghiệp, có đủ thống kê để tính cỡ ảnh hưởng
(Pearson \(r\), \(\beta\), \(t\)-stat).

\hypertarget{cuxf4ng-cux1ee5-truxedch-xuux1ea5t-dux1eef-liux1ec7u-vuxe0-cuxf4ng-bux1ed1-sux1eed-dux1ee5ng-truxed-tuux1ec7-nhuxe2n-tux1ea1o}{%
\subsection{3.2.4 Công cụ trích xuất dữ liệu và công bố sử dụng trí tuệ
nhân
tạo}\label{cuxf4ng-cux1ee5-truxedch-xuux1ea5t-dux1eef-liux1ec7u-vuxe0-cuxf4ng-bux1ed1-sux1eed-dux1ee5ng-truxed-tuux1ec7-nhuxe2n-tux1ea1o}}

Việc trích xuất cỡ ảnh hưởng từ K=288 ước lượng (k=238 nghiên cứu) được
thực hiện bằng phần mềm chuyên dụng \textbf{M-AIDA} (Meta-Analysis
Intelligent Data Assistant, phiên bản 7.0) do chính nghiên cứu sinh phát
triển và đã nộp hồ sơ đăng ký quyền tác giả. M-AIDA dùng mô hình ngôn
ngữ lớn (large language model, LLM) để đọc và trích các thống kê gốc (cỡ
mẫu \(N\), hệ số tương quan \(r\), thống kê \(t\), bậc tự do \(df\), hệ
số hồi quy chuẩn hóa \(\beta\), giá trị \(p\), khoảng tin cậy) từ bản
PDF của từng nghiên cứu, kèm điểm tin cậy trích xuất (extraction
confidence) ba mức --- 1,0 khi lấy trực tiếp \(r\); 0,8 khi suy từ
\(t\); 0,6 khi suy từ \(\beta\) theo Peterson và Brown (2005) --- và tự
động gắn cờ cần xác minh đối với mọi bản ghi có độ tin cậy dưới 0,7.

Để bảo đảm liêm chính và khả năng tái lập, luận án áp dụng nguyên tắc
\textbf{con người trong vòng lặp (human-in-the-loop) với kiểm chứng toàn
bộ (100\%) bởi nghiên cứu sinh} trong vai trò chủ trì nghiên cứu
(Principal Investigator): (i) LLM \textbf{chỉ} trích xuất thống kê và
\textbf{không} suy đoán các biến điều tiết thể chế --- các nhãn chế độ
ICRV, pha DPL và chỉ số cDAI do nghiên cứu sinh gán thủ công từ bảng tra
cứu ngoài (WGI Rule of Law, năm dữ liệu trung vị, Digital Adoption
Index), bảo đảm tính kiểm chứng độc lập; (ii) mỗi bản ghi phải được
nghiên cứu sinh phê duyệt --- với quyền ghi đè bất kỳ trường nào ---
trước khi được \textbf{khóa bất biến (irreversible data lock)} kèm dấu
thời gian; (iii) chỉ các bản ghi đã khóa mới được đưa vào tập phân tích
tổng hợp. Cơ chế hai bước phê duyệt--khóa tạo dấu vết kiểm toán (audit
trail) phân tách rạch ròi phần do máy trích xuất và phần do con người
quyết định, phù hợp với chuẩn mực công bố sử dụng trí tuệ nhân tạo ngày
càng được các tạp chí học thuật yêu cầu. Kiến trúc, quy trình trích
xuất--chuyển đổi cỡ ảnh hưởng và cơ chế quản trị dữ liệu của M-AIDA được
trình bày chi tiết ở \textbf{Phụ lục B}.

\hypertarget{giai-ux111oux1ea1n-2-phuxe2n-tuxedch-thux1ef1c-nghiux1ec7m-ux111a-quux1ed1c-gia}{%
\section{3.3 Giai đoạn 2: Phân tích thực nghiệm đa quốc
gia}\label{giai-ux111oux1ea1n-2-phuxe2n-tuxedch-thux1ef1c-nghiux1ec7m-ux111a-quux1ed1c-gia}}

\hypertarget{nguux1ed3n-dux1eef-liux1ec7u}{%
\subsection{3.3.1 Nguồn dữ liệu}\label{nguux1ed3n-dux1eef-liux1ec7u}}

Luận án sử dụng World Bank Enterprise Surveys cho \textbf{50 nền kinh tế
châu Á và Thái Bình Dương} trong giai đoạn \textbf{2006--2026} (88.869
doanh nghiệp trong khung phân tích, 103 cặp nền kinh tế và năm; bao gồm
Nhật Bản 2025; pool phân loại ICRV 96.415 doanh nghiệp/52 nhãn nền)
(World Bank, n.d., 2019, 2023, 2026). WBES là bộ dữ liệu doanh nghiệp
chuẩn hóa, phù hợp cho nghiên cứu so sánh đa quốc gia về cấp doanh
nghiệp kết quả (Aterido et al., 2011). Quy trình hợp nhất và hài hòa hóa
chi tiết các bộ khảo sát quốc gia--năm thành bộ dữ liệu phân tích thống
nhất (sáu bước S1--S6, xử lý tính so sánh tiền tệ, chiến lược lát cắt
ngang gộp với hiệu ứng cố định hai chiều, cùng dòng dữ liệu PRISMA và mã
tái lập) được trình bày đầy đủ ở \textbf{Phụ lục A}.

\textbf{Cơ sở kế thừa}: WBES đã được sử dụng rộng rãi trong nghiên cứu
IB về thị trường mới nổi (Ayyagari et al., 2011; Cuervo-Cazurra et al.,
2018; Chen \& Meng, 2022), đồng thời đã được sử dụng trong bài nghiên
cứu thực trạng châu Á mới nổi với 17 nền kinh tế và
\textasciitilde40.633 quan sát doanh nghiệp (nghiên cứu thành phần P1).

\textbf{Đóng góp mới}: mở rộng quy mô từ 17 sang \textbf{50 nền kinh tế
châu Á và Thái Bình Dương}, với 6 chế độ con ICRV (Advanced
innovation-driven, tiên tiến tài nguyên, trung bình cao, chuyển đổi thu
nhập trung bình thấp, đang nổi, SIDS). Quy mô này là lớn nhất cho khu
vực trong tài liệu quốc tế hóa--performance, gồm cả trường hợp biên đảo
nhỏ Thái Bình Dương (7 nước, n=959 (mẫu phân tích sau lọc missing)) cho
phép kiểm định forced quốc tế hóa penalty (nghiên cứu thành phần P8).

\begin{quote}
\textbf{Ghi chú về trọng số khảo sát (survey sampling weights):} WBES
cung cấp sampling weights để đại diện cho tổng thể doanh nghiệp trong
từng quốc gia. Tuy nhiên, luận án \emph{không áp dụng sampling weights}
trong ước lượng mô hình hồi quy vì ba lý do: (1) Mục tiêu của luận án là
kiểm định cơ chế lý thuyết (mechanism testing) chứ không phải ước lượng
mô tả đại diện tổng thể (population-representative descriptive
estimation), trong bối cảnh này, không áp dụng weights là thực hành
chuẩn trong tài liệu IB sử dụng WBES (Aterido et al., 2011;
Cuervo-Cazurra et al., 2018); (2) hiệu ứng cố định ở cấp cặp nền kinh tế
và năm đã kiểm soát đặc điểm cấu trúc của từng sóng khảo sát; (3) Việc
áp dụng weights có thể làm sai lệch ước lượng khi mẫu được gộp chung từ
nhiều quốc gia có quy mô tổng thể doanh nghiệp khác nhau nhiều bậc (ví
dụ: Trung Quốc \textasciitilde100 triệu so với Tonga \textasciitilde500
doanh nghiệp). Phân tích độ nhạy có áp dụng survey weights cho từng quốc
gia riêng lẻ (P3/P4/P5) không thay đổi hướng hoặc mức độ ý nghĩa của các
hệ số chính (xem Mục 3.5).
\end{quote}

\hypertarget{huxe0i-huxf2a-huxf3a-ba-thux1ebf-hux1ec7-lux1b0ux1ee3c-ux111ux1ed3-bux1ea3ng-hux1ecfi-wbes}{%
\subsection{3.3.1.1 Hài hòa hóa ba thế hệ lược đồ bảng hỏi
WBES}\label{huxe0i-huxf2a-huxf3a-ba-thux1ebf-hux1ec7-lux1b0ux1ee3c-ux111ux1ed3-bux1ea3ng-hux1ecfi-wbes}}

Bộ công cụ WBES đã trải qua ba thế hệ lược đồ với hệ thống mã biến và
phạm vi phân hệ khác nhau: thế hệ PICS3 đời đầu, thế hệ Standardized
toàn cầu áp dụng từ năm 2006, và thế hệ BREADY/BEE gần đây. Sự thay đổi
này tạo ra thách thức then chốt cho hợp nhất đa quốc gia, vì cùng một
khái niệm có thể được hỏi bằng trường khác nhau hoặc thiếu hẳn ở một số
đợt. Luận án giải quyết bằng một lược đồ biến chung ánh xạ mọi trường dị
biệt về các module cốt lõi bất biến của bộ công cụ toàn cầu, với tiêu
chí vận hành để một đợt được nhận vào khung so sánh là sự hiện diện của
các biến lõi đã hài hòa, cụ thể là doanh thu (d2), lao động (l1) và xuất
khẩu trực tiếp cùng gián tiếp (d3b và d3c). Đợt khảo sát đáp ứng năm
khảo sát từ 2006 nhưng chạy trên bộ công cụ khu vực cũ không chứa các
biến lõi này bị loại do bất tương thích bộ công cụ chứ không phải do
thiếu dữ liệu.

Quy trình hợp nhất tuân theo nguyên tắc minh bạch kiểu PRISMA cho dữ
liệu thứ cấp, gồm sáu bước tuần tự: lọc loại khảo sát, khử trùng lặp một
khảo sát cho mỗi cặp nền kinh tế và năm, hài hòa hóa biến, mã hóa mã
không phản hồi của WBES thành giá trị khuyết, phân tầng chế độ thể chế,
đến tạo mẫu phân tích bằng xóa theo danh sách. Mỗi bước ghi rõ tiêu chí
và số quan sát còn lại nhằm cho phép tái dựng và kiểm chứng. Một hệ quả
của xóa theo danh sách là cỡ mẫu giảm dần khi thêm biến kiểm soát, rõ
nhất ở biến sở hữu nước ngoài vốn khuyết nhiều ở các nền nhỏ và sóng
sớm; luận án báo cáo minh bạch sự suy giảm này và kiểm tra tính ổn định
của điểm uốn qua các mô hình để bảo đảm kết quả không bị chi phối bởi cơ
chế khuyết. Toàn bộ thao tác chi tiết của quy trình sáu bước, cùng dòng
dữ liệu kiểu PRISMA và mã tái lập, được trình bày đầy đủ ở Phụ lục A.

\hypertarget{ux111o-lux1b0ux1eddng-biux1ebfn-phux1ee5-thuux1ed9c-hiux1ec7u-quux1ea3-houx1ea1t-ux111ux1ed9ng}{%
\subsection{3.3.2 Đo lường biến phụ thuộc: Hiệu quả hoạt
động}\label{ux111o-lux1b0ux1eddng-biux1ebfn-phux1ee5-thuux1ed9c-hiux1ec7u-quux1ea3-houx1ea1t-ux111ux1ed9ng}}

\textbf{Biến chính}: năng suất lao động \(\mathrm{LP}\) (labour
productivity), đo bằng \(\ln(\mathrm{LP}) = \ln(d2/l1)\), trong đó
\(d2\) là tổng doanh thu năm tài chính gần nhất và \(l1\) là số lao động
thường xuyên toàn thời gian (mã trường WBES).

\textbf{Cơ sở kế thừa}: năng suất lao động được sử dụng rộng rãi trong
nghiên cứu WBES và trong tài liệu về năng suất doanh nghiệp (Bloom et
al., 2012; Hsieh \& Klenow, 2009). Chỉ tiêu này đã được vận hành hóa
thành công trong bài thực trạng châu Á mới nổi (nghiên cứu thành phần
P1).

\textbf{Lập luận chọn}: WBES không phải bộ báo cáo tài chính đầy đủ, nên
năng suất lao động phù hợp hơn các tỷ suất sinh lời kế toán (ROA, ROE,
ROS) trong bối cảnh đa quốc gia. Ngoài ra, năng suất lao động ít bị ảnh
hưởng bởi chuẩn mực kế toán khác biệt giữa các nước hơn so với các chỉ
tiêu kế toán (Combs et al., 2005; Richard et al., 2009).

\textbf{Nhất quán đo lường xuyên các nghiên cứu thành phần và điều kiện
so sánh tiền tệ}: cùng một cấu trúc \(\ln(\mathrm{LP})\) từ trường WBES
\(d2/l1\) được dùng thống nhất ở P3, P4, P5 và P7, nhưng cách xử lý đơn
vị tiền tệ khác nhau \textbf{theo đúng bản chất thiết kế} của từng
nghiên cứu: (i) các nghiên cứu \textbf{đơn quốc gia} (P3 Việt Nam, P4
Singapore, P5 Trung Quốc, P9 Ấn Độ) đo doanh thu bằng \textbf{một nội tệ
duy nhất} nên mức \(\ln(\mathrm{LP})\) so sánh được trực tiếp trong nội
bộ mẫu; khác biệt mặt bằng giá giữa các đợt khảo sát được hấp thụ bằng
biến giả đợt, hiệu ứng cố định năm (P5: \(D_{2024}\); P3 bổ sung hiệu
chỉnh PPP cho khoảng cách 14 năm 2009--2023), kèm cắt đuôi 1\%/99\%
(P4); (ii) nghiên cứu \textbf{đa quốc gia} (P7) gộp doanh thu bằng
\textbf{nhiều nội tệ}, do đó mức \(\ln(\mathrm{LP})\) thô không so sánh
được giữa các nước, tính so sánh được bảo đảm bằng chuẩn hóa z within
country--year và hiệu ứng cố định hai chiều, như trình bày chi tiết ở
Phụ lục A (Mục A.5); mọi so sánh \textbf{mức} xuyên quốc gia trong luận
án dùng tỷ suất không thứ nguyên (ROS) hoặc giá trị đã hiệu chỉnh PPP.
Như vậy tính so sánh xuyên quốc gia của thước đo là \textbf{có điều
kiện}, đạt được qua thiết kế ước lượng, không phải tự động ở mức thô. Về
phía biến quốc tế hóa, khung ước lượng đa quốc gia của luận án (P7, P8,
P9) và thống kê mô tả ở Chuyên đề 1 đo \(\mathrm{FSTS}\) bằng
\textbf{tổng xuất khẩu trực tiếp và gián tiếp} (trường WBES
\(d3b + d3c\)), định nghĩa tham gia xuất khẩu chuẩn của bộ công cụ WBES;
hai nghiên cứu đơn quốc gia (P3 Việt Nam, P5 Trung Quốc) dùng
\textbf{riêng xuất khẩu trực tiếp} (\(d3c\)), vì xuất khẩu gián tiếp ủy
thác chức năng quốc tế hóa cho trung gian thương mại với cơ chế học hỏi
và yêu cầu nguồn lực khác về bản chất (Hessels \& Terjesen, 2010; Peng
\& Ilinitch, 1998), trong khi tỷ trọng doanh thu xuất khẩu là thước đo
cường độ quốc tế hóa chuẩn của tài liệu về quan hệ quốc tế hóa và hiệu
quả (Lu \& Beamish, 2001; Sullivan, 1994); lựa chọn thao tác hóa được
ghi rõ trong phần dữ liệu của từng nghiên cứu thành phần.

\textbf{Biến độ vững}: tỷ suất lợi nhuận trên doanh thu
(\(\mathrm{ROS}\)), tăng trưởng doanh thu, tăng trưởng lao động, và
logarit doanh thu năm \(\ln(d2)\).

\textbf{Đóng góp mới}: lập luận rõ ràng về hiệu quả hoạt động kinh doanh
của doanh nghiệp là khái niệm đa chiều và việc lựa chọn thước đo phải
phù hợp với bối cảnh dữ liệu (xem Chương 2).

\hypertarget{vux1eadn-huxe0nh-huxf3a-thux1b0ux1edbc-ux111o-nux103ng-suux1ea5t-vuxe0-quy-truxecnh-xux1eed-luxfd-phuxe2n-phux1ed1i}{%
\subsection{3.3.2.1 Vận hành hóa thước đo năng suất và quy trình xử lý
phân
phối}\label{vux1eadn-huxe0nh-huxf3a-thux1b0ux1edbc-ux111o-nux103ng-suux1ea5t-vuxe0-quy-truxecnh-xux1eed-luxfd-phuxe2n-phux1ed1i}}

Năng suất lao động được dựng từ hai trường khảo sát cốt lõi của bộ công
cụ WBES: tổng doanh thu năm tài chính gần nhất (trường d2) và số lao
động thường xuyên toàn thời gian cuối kỳ (trường l1). Biến phụ thuộc
nhận dạng logarit tự nhiên của thương số hai trường này, nghĩa là
\(\ln(\mathrm{LP}) = \ln(d2 / l1)\), một phép biến đổi vừa nén đuôi phải
của phân phối doanh thu vừa cho phép diễn giải hệ số theo phần trăm. Do
doanh thu thô được báo cáo bằng nội tệ của từng nền kinh tế, mức năng
suất thô không so sánh trực tiếp giữa các nước, một mối đe dọa hiệu lực
được xử lý qua thiết kế ước lượng chứ không qua quy đổi tỷ giá danh
nghĩa.

Trước khi ước lượng, năng suất lao động trải qua hai bước xử lý phân
phối nhằm bảo đảm tính bền vững của ước lượng trước quan sát cực trị.
Bước thứ nhất là cắt đuôi (winsorize) tại ngưỡng phân vị 1 và 99 áp dụng
trong từng cặp nền kinh tế và năm, qua đó giới hạn ảnh hưởng của giá trị
ngoại lai mà không loại bỏ quan sát khỏi mẫu. Bước thứ hai, chỉ áp dụng
cho khung ước lượng đa quốc gia, là chuẩn hóa z trong từng cặp nền kinh
tế và năm, đưa mọi quan sát về cùng thang vị thế tương đối trong phân
phối nội bộ nước--năm. Phép chuẩn hóa này khử triệt để khác biệt đơn vị
tiền tệ và mặt bằng giá, đồng thời bảo toàn cấu trúc biến thiên bên
trong mỗi nước, nơi diễn ra việc nhận dạng quan hệ quốc tế hóa--hiệu
quả. Đối với các nghiên cứu đơn quốc gia, doanh thu được biểu thị bằng
một nội tệ duy nhất nên mức \(\ln(\mathrm{LP})\) so sánh được trực tiếp
trong nội bộ mẫu; khác biệt mặt bằng giá giữa các đợt khảo sát được hấp
thụ bằng biến giả đợt hoặc hiệu ứng cố định năm, kèm cắt đuôi 1 và 99.

\hypertarget{ux111o-lux1b0ux1eddng-biux1ebfn-ux111ux1ed9c-lux1eadp-mux1ee9c-ux111ux1ed9-quux1ed1c-tux1ebf-huxf3a}{%
\subsection{3.3.3 Đo lường biến độc lập: Mức độ quốc tế
hóa}\label{ux111o-lux1b0ux1eddng-biux1ebfn-ux111ux1ed9c-lux1eadp-mux1ee9c-ux111ux1ed9-quux1ed1c-tux1ebf-huxf3a}}

\textbf{Khái niệm.} Biến độc lập trọng tâm của luận án là \emph{mức độ
quốc tế hóa} (degree of internationalization, DOI) --- phạm vi mà doanh
nghiệp gắn hoạt động của mình với thị trường nước ngoài. Theo Sullivan
(1994), mức độ quốc tế hóa là khái niệm đa chiều gồm ba thành phần:
chiều theo doanh số (tỷ trọng doanh thu nước ngoài), chiều theo tài sản
(tỷ trọng tài sản đặt ở nước ngoài) và chiều theo nhân sự, tâm lý (mức
độ phân tán quốc tế của nhân lực và định hướng quản trị). Bộ dữ liệu
WBES chỉ thu thập nhất quán \textbf{chiều theo doanh số} xuyên ba thế hệ
bảng hỏi, nên luận án thao tác hóa mức độ quốc tế hóa bằng \emph{cường
độ xuất khẩu} \(\mathrm{FSTS}\) (foreign sales to total sales) --- tỷ
trọng doanh thu xuất khẩu trên tổng doanh thu, miền giá trị \([0,1]\).

\textbf{Biến chính.} \(\mathrm{FSTS}\) là tỷ lệ doanh thu xuất khẩu trên
tổng doanh thu. Mô hình đưa vào đồng thời số hạng bậc nhất
\(\mathrm{FSTS}\) và số hạng bậc hai \(\mathrm{FSTS}^{2}\) để kiểm định
quan hệ phi tuyến dạng chữ U ngược.

\textbf{Cơ sở kế thừa}: \(\mathrm{FSTS}\) là thước đo đơn chiều của mức
độ quốc tế hóa được dùng phổ biến nhất trong tài liệu về quan hệ quốc tế
hóa và hiệu quả khi dữ liệu chiều tài sản và chiều nhân sự quốc tế không
khả dụng (Sullivan, 1994; Lu \& Beamish, 2001; Hitt et al., 2006). Việc
bổ sung \(\mathrm{FSTS}^{2}\) để kiểm định phi tuyến đã được áp dụng
trong Lu và Beamish (2004), Contractor et al.~(2003), và Hitt et
al.~(1997). Dạng bậc hai được nghiên cứu thành phần P2 xác nhận có ý
nghĩa thống kê ở Trung Quốc với điểm uốn khoảng 47,8\%
\(\mathrm{FSTS}\), làm cơ sở cho việc kiểm định mở rộng trong luận án.
Giới hạn đo lường đơn chiều này (chỉ chiều doanh số) được nêu minh bạch
như một hạn chế nghiên cứu ở Chương 5.

\hypertarget{vux1eadn-huxe0nh-huxf3a-cux1b0ux1eddng-ux111ux1ed9-quux1ed1c-tux1ebf-huxf3a-vuxe0-quy-ux1b0ux1edbc-trung-tuxe2m-huxf3a}{%
\subsection{3.3.3.1 Vận hành hóa cường độ quốc tế hóa và quy ước trung
tâm
hóa}\label{vux1eadn-huxe0nh-huxf3a-cux1b0ux1eddng-ux111ux1ed9-quux1ed1c-tux1ebf-huxf3a-vuxe0-quy-ux1b0ux1edbc-trung-tuxe2m-huxf3a}}

Cường độ quốc tế hóa được đo bằng tỷ trọng doanh thu xuất khẩu trên tổng
doanh thu (Foreign Sales To Total Sales, FSTS), giới hạn miền giá trị từ
0 đến 1. Doanh nghiệp không xuất khẩu nhận giá trị FSTS bằng 0, một giá
trị hợp lệ phản ánh trạng thái phi xuất khẩu chứ không phải dữ liệu
khuyết. Khung đa quốc gia cộng gộp xuất khẩu trực tiếp và gián tiếp
(tổng hai trường d3b và d3c), đúng định nghĩa tham gia xuất khẩu chuẩn
của bộ công cụ WBES; hai nghiên cứu đơn quốc gia (Việt Nam, Trung Quốc)
dùng riêng xuất khẩu trực tiếp (trường d3c), vì xuất khẩu gián tiếp ủy
thác chức năng quốc tế hóa cho trung gian thương mại với cơ chế học hỏi
và yêu cầu nguồn lực khác về bản chất.

Để kiểm định phi tuyến, mô hình đưa đồng thời số hạng bậc nhất và số
hạng bậc hai của cường độ xuất khẩu. Trước khi tạo số hạng bậc hai,
cường độ xuất khẩu được trung tâm hóa quanh trung bình trong từng đợt
khảo sát hoặc trong từng mẫu gộp, một quy ước làm giảm tương quan giả
tạo giữa số hạng tuyến tính và số hạng bình phương, qua đó cải thiện độ
ổn định số học của ước lượng và làm rõ diễn giải hệ số tại điểm trung
bình của mẫu. Điểm uốn của quan hệ chữ U ngược được khôi phục về thang
gốc bằng cách cộng trung bình mẫu vào tỷ số
\(-\hat\beta_1 / (2\hat\beta_2)\), cho phép báo cáo ngưỡng tối ưu theo
tỷ trọng xuất khẩu thực tế. Dạng bậc hai được dùng nhất quán trong mọi
nghiên cứu thành phần, vì nó đủ để nhận dạng một điểm uốn nội miền và
phù hợp với mật độ phân phối cường độ xuất khẩu của các mẫu doanh nghiệp
WBES; luận án không mở rộng lên dạng bậc ba để tránh nhận dạng điểm uốn
không ổn định ở vùng cường độ xuất khẩu cao vốn thưa quan sát.

\hypertarget{vux1eadn-huxe0nh-huxf3a-chi-tiux1ebft-cux1b0ux1eddng-ux111ux1ed9-xuux1ea5t-khux1ea9u-xuyuxean-cuxe1c-nghiuxean-cux1ee9u-thuxe0nh-phux1ea7n}{%
\subsection{3.3.3.2 Vận hành hóa chi tiết cường độ xuất khẩu xuyên các
nghiên cứu thành
phần}\label{vux1eadn-huxe0nh-huxf3a-chi-tiux1ebft-cux1b0ux1eddng-ux111ux1ed9-xuux1ea5t-khux1ea9u-xuyuxean-cuxe1c-nghiuxean-cux1ee9u-thuxe0nh-phux1ea7n}}

Sự thống nhất khái niệm của cường độ xuất khẩu xuyên các nghiên cứu
thành phần không kéo theo một định nghĩa thao tác duy nhất, mà là một họ
định nghĩa được hiệu chỉnh theo bản chất bộ công cụ của từng đợt và theo
câu hỏi nghiên cứu. Khung đa quốc gia cộng gộp xuất khẩu trực tiếp và
gián tiếp rồi chia cho một trăm để đưa về tỷ trọng, cụ thể
\(\mathrm{FSTS} = (d3b + d3c)/100\) với miền giá trị từ 0 đến 1, đúng
định nghĩa tham gia xuất khẩu chuẩn của bộ công cụ toàn cầu; trung bình
mẫu của tỷ trọng này trong mẫu gộp năm mươi nền là 9,0\%. Ngược lại, hai
nghiên cứu đơn quốc gia dùng riêng xuất khẩu trực tiếp dựng từ một
trường khảo sát: ở Trung Quốc cường độ xuất khẩu là tỷ trọng doanh thu
xuất khẩu trực tiếp, còn ở đảo nhỏ Thái Bình Dương cường độ xuất khẩu
được dựng từ trường \texttt{d3c} chia cho một trăm. Lý do lựa chọn riêng
xuất khẩu trực tiếp đã nêu ở Mục 3.3.3.1: xuất khẩu gián tiếp ủy thác
chức năng quốc tế hóa cho trung gian thương mại, mang cơ chế học hỏi và
yêu cầu nguồn lực khác về bản chất.

Doanh nghiệp không xuất khẩu nhận giá trị cường độ bằng 0 và được giữ
trong mẫu, vì giá trị này phản ánh trạng thái phi xuất khẩu hợp lệ chứ
không phải dữ liệu khuyết; điểm uốn của quan hệ do đó được nhận dạng
trên toàn mẫu chứ không chỉ trên tập doanh nghiệp xuất khẩu. Quy ước này
đặc biệt hệ trọng ở đảo nhỏ Thái Bình Dương, nơi chỉ 132 trong số 959
quan sát phân tích, tương đương 13,8\%, là doanh nghiệp có xuất khẩu với
cường độ dương: loại bỏ nhóm phi xuất khẩu sẽ làm sụp đổ công suất thống
kê và làm sai lệch ước lượng ở biên tham gia.

Quy ước trung tâm hóa được áp dụng nhất quán nhưng theo hai cấp khác
nhau. Ở mẫu gộp đa quốc gia, cường độ xuất khẩu được trung tâm hóa quanh
trung bình toàn mẫu trước khi lập bình phương. Ở các nghiên cứu đa đợt,
cường độ xuất khẩu được trung tâm hóa trong từng đợt nhằm chặn khác biệt
thành phần mẫu giữa các đợt làm nhiễu vị trí điểm uốn. Trung tâm hóa làm
giảm tương quan giả tạo giữa số hạng tuyến tính và số hạng bình phương
mà không làm dịch chuyển vị trí điểm uốn, vì công thức điểm uốn
\(-\hat\beta_1/(2\hat\beta_2)\) bất biến với phép trung tâm hóa; điểm
uốn trên thang trung tâm hóa được quy đổi về thang gốc bằng cách cộng
trung bình mẫu, tức
\(\mathrm{TP} = -\hat\beta_1/(2\hat\beta_2) + \overline{\mathrm{FSTS}}\).

\hypertarget{ux111o-lux1b0ux1eddng-biux1ebfn-ux111iux1ec1u-tiux1ebft}{%
\subsection{3.3.4 Đo lường biến điều
tiết}\label{ux111o-lux1b0ux1eddng-biux1ebfn-ux111iux1ec1u-tiux1ebft}}

\hypertarget{cuxe1c-cux1ea5u-truxfac-sux1ed1-nux103ng-lux1ef1c-cuxf4ng-nghux1ec7-tci-vuxe0-chux1ea5p-nhux1eadn-sux1ed1-dai}{%
\subsubsection{Các cấu trúc số: năng lực công nghệ (TCI) và chấp nhận số
(DAI)}\label{cuxe1c-cux1ea5u-truxfac-sux1ed1-nux103ng-lux1ef1c-cuxf4ng-nghux1ec7-tci-vuxe0-chux1ea5p-nhux1eadn-sux1ed1-dai}}

Luận án phân biệt rõ hai cấu trúc:

\begin{itemize}
\tightlist
\item
  \textbf{Chỉ số năng lực công nghệ} (technological capability index,
  \(\mathrm{TCI}\)): chiều sâu năng lực công nghệ, tổng hợp từ cấp phép
  công nghệ nước ngoài, chi tiêu nghiên cứu và phát triển, đổi mới sản
  phẩm, và chứng nhận chất lượng.
\item
  \textbf{Chỉ số chấp nhận số} (digital adoption index,
  \(\mathrm{DAI}\)): mức độ chấp nhận số bề mặt, tổng hợp từ website,
  thư điện tử, bán hàng trực tuyến (\(\mathrm{DAI}_{\text{thin}}\)) hoặc
  bổ sung thanh toán điện tử và nhà cung cấp điện tử
  (\(\mathrm{DAI}_{\text{rich}}\)).
\end{itemize}

\textbf{Nguyên tắc không chồng lấn}: TCI và DAI không chia sẻ thành phần
chỉ báo để bảo đảm độ thuần cấu trúc. Đặc biệt, cấp phép công nghệ nước
ngoài thuộc TCI chứ không được đồng thời tính vào DAI.

\textbf{Biến thể đo lường (rút gọn/đầy đủ).} Do bộ công cụ WBES thay đổi
qua ba thế hệ bảng hỏi (PICS3, Standardized, BREADY) khiến không phải
mọi mục đều khả dụng ở mọi sóng, TCI được vận hành hóa thống nhất theo
hai biến thể của cùng một chỉ số hợp thành cấu thành:

\begin{itemize}
\tightlist
\item
  \(\mathrm{TCI}_{\text{thin}}\) = chuẩn hóa z của trung bình hai mục
  cốt lõi: chứng chỉ chất lượng quốc tế (\texttt{b8}) và cấp phép công
  nghệ nước ngoài (\texttt{e6}). Dùng trong các nghiên cứu Việt Nam
  (P3), Singapore (P4), đa quốc gia (P7) và đảo nhỏ Thái Bình Dương
  (P8).
\item
  \(\mathrm{TCI}_{\text{full}}\) = chuẩn hóa z của trung bình bốn mục:
  \texttt{b8}, \texttt{e6}, đổi mới sản phẩm (\texttt{h1}), và chi
  nghiên cứu--phát triển (\texttt{h8}). Dùng trong nghiên cứu Trung Quốc
  (P5), nơi cả hai sóng 2012 và 2024 đều có đủ phân hệ đổi mới.
\end{itemize}

Bộ mục chính xác của mỗi nghiên cứu thành phần được nêu lại trong bảng
định nghĩa biến của nghiên cứu đó (Mục 3.4.5) và Phụ lục A; khác biệt
giữa biến thể rút gọn và đầy đủ phản ánh ràng buộc khả dụng dữ liệu
khách quan, không phải lựa chọn tùy ý. Vì các mục là chỉ báo nhị phân
của cùng một cấu trúc cấu thành (năng lực công nghệ), hai biến thể tương
quan cao và cho kết luận điều tiết nhất quán về dấu; phần độ vững (Mục
3.5.1) kiểm tra độ nhạy của kết quả với lựa chọn biến thể. Bộ mục TCI
của nghiên cứu đa quốc gia (P7) đã được đối chiếu trực tiếp với mã xây
dựng dữ liệu thực tế (quy trình xử lý dữ liệu thô WBES 50 nền kinh tế,
gồm Nhật Bản 2025): biến \texttt{tci\_z} được dựng đúng từ hai mục b8 và
e6, thống nhất với khung rút gọn nêu trên.

\textbf{Cơ sở kế thừa}: Bhandari et al.~(2023) phân biệt số hóa, quốc tế
hóa và năng lực theo logic điều phối nguồn lực; Verhoef et al.~(2021)
phân biệt ba giai đoạn (số hóa dữ liệu, số hóa quy trình, chuyển đổi
số). Khuôn mẫu ``hiệu ứng lá chắn số'' được nghiên cứu thành phần P1
phát hiện trên 17 nước châu Á mới nổi.

\textbf{Đóng góp mới}: Đưa ra \textbf{giao thức hài hòa hóa đo lường}
trong bối cảnh WBES, đảm bảo độ thuần cấu trúc ở 50 nền kinh tế châu Á
và Thái Bình Dương, với quy trình hòa hợp 3 thế hệ lược đồ. Phân tách
TCI và DAI theo (i) Bharadwaj et al.~(2013), năng lực công nghệ thông
tin so với chiến lược kinh doanh số có mạng quan hệ lý thuyết khác nhau;
(ii) phân tầng 3 cấp của Verhoef et al.~(2021), số hóa dữ liệu, số hóa
quy trình, chuyển đổi số; (iii) truyền thống Lall (1992) và Cohen \&
Levinthal (1990) coi năng lực công nghệ là chiều sâu năng lực hấp thụ
nội tại (absorptive capacity), khác về bản chất với chấp nhận số ở giao
diện ngoại tại; (iv) logic điều phối nguồn lực của Bhandari et
al.~(2023) cho quan hệ quốc tế hóa và hiệu quả; và (v) 4 tiêu chí chỉ số
hợp thành cấu thành của Coltman et al.~(2008), đảm bảo TCI và DAI là hai
chỉ số hợp thành độc lập.

\hypertarget{biux1ebfn-thux1ec3-chux1ebf}{%
\subsubsection{Biến thể chế}\label{biux1ebfn-thux1ec3-chux1ebf}}

\textbf{Biến chính}: chỉ số rào cản kinh doanh (business obstacle index,
tổng hợp từ các câu hỏi WBES về rào cản hoạt động), các biến đại diện
chất lượng quản trị cấp quốc gia, chỉ số pháp quyền WGI (Kaufmann et
al., 2011), và 6 chế độ con ICRV (phân loại mở rộng từ Khanna \& Palepu,
2010).

\textbf{Cơ sở kế thừa}: Khanna và Palepu (2010) với khoảng trống thể
chế; North (1990) với phân tích thể chế; Marano et al.~(2016) với tổng
hợp bằng chứng về các biến điều tiết thể chế.

\hypertarget{ux111ux1eb7c-ux111iux1ec3m-nhuxe0-quux1ea3n-trux1ecb-cux1ea5p-cao}{%
\subsubsection{Đặc điểm nhà quản trị cấp
cao}\label{ux111ux1eb7c-ux111iux1ec3m-nhuxe0-quux1ea3n-trux1ecb-cux1ea5p-cao}}

\textbf{Biến chính}: kinh nghiệm nhà quản trị cấp cao (số năm kinh
nghiệm trong ngành) và giới tính nhà quản trị cấp cao (nữ/nam).

\textbf{Cơ sở kế thừa}: Hambrick và Mason (1984), Hambrick (2007), Hsu
et al.~(2013), Post và Byron (2015).

\textbf{Điều kiện sử dụng}: chỉ trong tập con có đủ dữ liệu; nếu thiếu,
đưa vào phần kiểm định độ vững chứ không bắt buộc trong mô hình chính.

\hypertarget{phuxe2n-tux1ea7ng-chux1ec9-sux1ed1-chux1ea5p-nhux1eadn-sux1ed1-vuxe0-quy-tux1eafc-xuxe2y-dux1ef1ng-chux1ec9-sux1ed1-hux1ee3p-thuxe0nh-cux1ea5u-thuxe0nh}{%
\subsection{3.3.4.1 Phân tầng chỉ số chấp nhận số và quy tắc xây dựng
chỉ số hợp thành cấu
thành}\label{phuxe2n-tux1ea7ng-chux1ec9-sux1ed1-chux1ea5p-nhux1eadn-sux1ed1-vuxe0-quy-tux1eafc-xuxe2y-dux1ef1ng-chux1ec9-sux1ed1-hux1ee3p-thuxe0nh-cux1ea5u-thuxe0nh}}

Chỉ số chấp nhận số (Digital Adoption Index, DAI) được vận hành hóa theo
nguyên tắc phân tầng nhằm phản ánh chiều sâu tăng dần của số hóa doanh
nghiệp. Tầng 1 ghi nhận hiện diện số bề mặt qua chỉ báo website (trường
c22b), một mục nhị phân khả dụng ở hầu hết các đợt khảo sát. Tầng 2 bổ
sung các chức năng số kích hoạt giao dịch, cụ thể là thanh toán điện tử
và tỷ trọng doanh thu qua kênh điện tử (trường k33), những mục chỉ xuất
hiện ở các đợt khảo sát gần đây. Việc phân tầng tạo ra ba biến thể đo
lường: biến thể chỉ Tầng 1 (đơn biến nhị phân) dùng ở Việt Nam và đảo
nhỏ Thái Bình Dương, biến thể Tầng 1 cộng Tầng 2 dùng ở Singapore và
khung đa quốc gia khi trường k33 có mặt, và biến thể thu gọn về Tầng 1
khi trường Tầng 2 thiếu. Khác biệt giữa các biến thể phản ánh ràng buộc
khả dụng dữ liệu khách quan qua ba thế hệ bảng hỏi chứ không phải lựa
chọn tùy ý.

Cả chỉ số năng lực công nghệ (Technological Capability Index, TCI) và
DAI được xử lý như chỉ số hợp thành cấu thành (formative composite) chứ
không phải chỉ số phản ánh, vì mỗi chỉ báo nhị phân định nghĩa một khía
cạnh độc lập của khái niệm chứ không phải là biểu hiện tương quan của
một biến tiềm ẩn chung. Quy tắc xây dựng gồm ba bước: chuyển mỗi chỉ báo
thành dạng nhị phân 0--1, lấy trung bình cộng các chỉ báo trong cùng một
tầng, rồi chuẩn hóa z giá trị trung bình trong từng đợt khảo sát. Nguyên
tắc không chồng lấn được áp đặt nghiêm ngặt: cấp phép công nghệ nước
ngoài thuộc TCI và không được đồng thời tính vào DAI, bảo đảm hai chỉ số
là hai cấu trúc độc lập với mạng quan hệ lý thuyết khác nhau. Vì các chỉ
báo là biểu hiện nhị phân của cùng một cấu trúc cấu thành, các biến thể
của một chỉ số tương quan cao và cho kết luận điều tiết nhất quán về
dấu; phần độ vững kiểm tra trực tiếp độ nhạy của suy luận với lựa chọn
biến thể.

\hypertarget{uxe1nh-xux1ea1-mux1ee5c-khux1ea3o-suxe1t-cux1ee7a-chux1ec9-sux1ed1-nux103ng-lux1ef1c-cuxf4ng-nghux1ec7-vuxe0-chux1ec9-sux1ed1-chux1ea5p-nhux1eadn-sux1ed1}{%
\subsection{3.3.4.2 Ánh xạ mục khảo sát của chỉ số năng lực công nghệ và
chỉ số chấp nhận
số}\label{uxe1nh-xux1ea1-mux1ee5c-khux1ea3o-suxe1t-cux1ee7a-chux1ec9-sux1ed1-nux103ng-lux1ef1c-cuxf4ng-nghux1ec7-vuxe0-chux1ec9-sux1ed1-chux1ea5p-nhux1eadn-sux1ed1}}

Để bảo đảm tính minh bạch và tái lập, mỗi biến thể của hai chỉ số hợp
thành được truy nguyên về đúng bộ mục khảo sát WBES dùng để dựng nó.
Biến thể mỏng của chỉ số năng lực công nghệ được dựng từ đúng hai mục
nhị phân: chứng chỉ chất lượng được công nhận quốc tế (trường
\texttt{b8}) và cấp phép công nghệ nước ngoài (trường \texttt{e6}); giá
trị chỉ số là chuẩn hóa z của trung bình hai chỉ báo nhị phân này, xử lý
như một chỉ số hợp thành cấu thành. Bộ mục hai trường này được dùng nhất
quán ở khung đa quốc gia, ở Việt Nam, Singapore và đảo nhỏ Thái Bình
Dương. Biến thể đầy đủ chỉ dùng ở Trung Quốc, nơi cả hai đợt đều có đủ
phân hệ đổi mới, mở rộng lên bốn mục để nắm bắt chiều sâu năng lực nội
tại phong phú hơn so với biến thể hai mục dùng ở các nghiên cứu khác.
Khác biệt giữa biến thể mỏng và đầy đủ phản ánh ràng buộc khả dụng dữ
liệu khách quan qua ba thế hệ bảng hỏi chứ không phải lựa chọn tùy ý; vì
các mục là chỉ báo nhị phân của cùng một cấu trúc cấu thành, hai biến
thể cho kết luận điều tiết nhất quán về dấu, và phần độ vững (Mục 3.5.1)
kiểm tra trực tiếp độ nhạy với lựa chọn biến thể.

Chỉ số chấp nhận số được vận hành hóa theo nguyên tắc phân tầng đã trình
bày ở Mục 3.3.4.1, với bộ mục cụ thể cũng được truy nguyên về trường
khảo sát. Biến thể Tầng 1 dùng đúng một chỉ báo nhị phân là hiện diện
website vận hành (trường \texttt{c22b}); đây là biến thể duy nhất khả
dụng ở đảo nhỏ Thái Bình Dương, nơi bộ công cụ không thu thập các mục
Tầng 2 thanh toán điện tử và Tầng 3 đám mây/hệ hoạch định nguồn lực
doanh nghiệp. Khung đa quốc gia dùng biến thể Tầng 1 cộng Tầng 2 khi
trường thanh toán điện tử có mặt và thu gọn về Tầng 1 khi thiếu, đúng
như nguyên tắc nhúng tầng theo khả dụng bộ công cụ. Nguyên tắc không
chồng lấn được áp đặt nghiêm ngặt: cấp phép công nghệ nước ngoài thuộc
chỉ số năng lực công nghệ và không bao giờ được tính đồng thời vào chỉ
số chấp nhận số, bảo đảm hai chỉ số là hai cấu trúc độc lập với mạng
quan hệ lý thuyết khác nhau. Bộ mục năng lực công nghệ của khung đa quốc
gia đã được đối chiếu trực tiếp với mã xây dựng dữ liệu thực tế trên kho
năm mươi nền, gồm Nhật Bản 2025, xác nhận biến \texttt{tci\_z} được dựng
đúng từ hai mục \texttt{b8} và \texttt{e6}.

\hypertarget{biux1ebfn-kiux1ec3m-souxe1t}{%
\subsection{3.3.5 Biến kiểm soát}\label{biux1ebfn-kiux1ec3m-souxe1t}}

Khung thực nghiệm đưa vào một tập biến kiểm soát chuẩn trong nghiên cứu
dữ liệu doanh nghiệp và kinh doanh quốc tế, nhằm tách hiệu ứng của quốc
tế hóa và năng lực khỏi các yếu tố cấu trúc đồng biến với cả cường độ
xuất khẩu lẫn năng suất. Quy mô doanh nghiệp được đo bằng logarit tự
nhiên của số lao động toàn thời gian (mục l1 của WBES), một biến kiểm
soát thiết yếu vì doanh nghiệp lớn vừa có xác suất xuất khẩu cao hơn vừa
đạt năng suất cao hơn nhờ lợi thế quy mô; lấy logarit để xử lý phân phối
lệch phải đặc trưng của biến quy mô. Tuổi doanh nghiệp được tính bằng số
năm từ năm thành lập đến năm khảo sát, kiểm soát hiệu ứng học hỏi tích
lũy và hiệu ứng sống sót, theo đó doanh nghiệp tồn tại lâu hơn có thể đã
vượt qua các rào cản gia nhập và tích lũy năng lực.

Sở hữu nước ngoài được đưa vào dưới dạng biến giả nhận giá trị một khi
tỷ lệ sở hữu của nhà đầu tư nước ngoài đạt từ 10\% trở lên, ngưỡng chuẩn
của đầu tư trực tiếp nước ngoài, vì doanh nghiệp có vốn nước ngoài
thường tiếp cận mạng lưới quốc tế, công nghệ và thị trường theo cách
khác biệt về chất so với doanh nghiệp trong nước. Ngành được kiểm soát
bằng hệ biến giả ngành nhằm hấp thụ chênh lệch năng suất và cường độ
thương mại có hệ thống giữa các lĩnh vực. Quan trọng nhất, mọi mô hình
gộp đều mang hiệu ứng cố định nền kinh tế và hiệu ứng cố định năm: hiệu
ứng cố định nền kinh tế hấp thụ mọi khác biệt bất biến theo thời gian
giữa các nền kinh tế, gồm mặt bằng giá, chế độ tỷ giá, chất lượng thể
chế nền và quy mô thị trường nội địa; hiệu ứng cố định năm hấp thụ các
cú sốc vĩ mô chung như chu kỳ kinh tế toàn cầu và biến động thương mại
theo năm. Cấu trúc hiệu ứng cố định hai chiều này là then chốt vì nó
loại bỏ chính các nhiễu cấp nền kinh tế và cấp thời gian dễ tạo tương
quan giả giữa quốc tế hóa và hiệu quả, qua đó cô lập biến thiên trong
nội bộ nền kinh tế và năm để nhận dạng quan hệ trọng tâm.

\textbf{Cơ sở kế thừa}: tập biến kiểm soát và cấu trúc hiệu ứng cố định
tuân theo chuẩn trong nghiên cứu dữ liệu doanh nghiệp WBES và kinh doanh
quốc tế (Aterido et al., 2011; Ayyagari et al., 2011).

\hypertarget{vux1eadn-huxe0nh-huxf3a-chux1ebf-ux111ux1ed9-thux1ec3-chux1ebf-icrv-vuxe0-ux111ux1eb7c-ux111iux1ec3m-nhuxe0-quux1ea3n-trux1ecb}{%
\subsection{3.3.5.1 Vận hành hóa chế độ thể chế ICRV và đặc điểm nhà
quản
trị}\label{vux1eadn-huxe0nh-huxf3a-chux1ebf-ux111ux1ed9-thux1ec3-chux1ebf-icrv-vuxe0-ux111ux1eb7c-ux111iux1ec3m-nhuxe0-quux1ea3n-trux1ecb}}

Chế độ thể chế được vận hành hóa thành một biến phân loại số nguyên nhận
giá trị từ một đến sáu theo phân loại ICRV, đi từ chế độ tiên tiến đổi
mới đến chế độ đảo nhỏ Thái Bình Dương, trật tự tăng dần của số nhóm
tương ứng với chất lượng thể chế giảm dần. Biến này được dùng theo hai
vai trò trong khung đa quốc gia: làm biến điều tiết liên tục trong mô
hình tương tác, nơi hiệu ứng chính của nó bị hấp thụ bởi hiệu ứng cố
định nền kinh tế, và làm biến phân nhóm cho phân tích điểm uốn theo từng
chế độ. Lựa chọn xử lý ICRV như biến số nguyên liên tục thay vì hệ biến
giả là có chủ đích, nhằm giữ cỡ mẫu đủ lớn trong các mô hình điều tiết
và ba chiều, nơi việc tách thành sáu biến giả tương tác sẽ làm phân mảnh
mẫu và làm bất ổn ước lượng.

Phân tích điểm uốn theo từng chế độ làm lộ rõ cấu trúc ba vùng đặc trưng
của khu vực. Ở chế độ chuyển đổi thu nhập trung bình--thấp, vốn là chế
độ neo ước lượng gộp, quan hệ chữ U ngược định hình sắc nét với điểm uốn
quanh 43\%. Ở chế độ tiên tiến đổi mới, gồm cả Nhật Bản, điểm uốn nằm
gần hoặc vượt biên dữ liệu quan sát nên trong vùng dữ liệu quan hệ gần
như tuyến tính dương, với điểm uốn xấp xỉ 80\%. Ở các chế độ yếu nhất,
chữ U ngược mất cấu trúc; riêng ở chế độ đảo nhỏ Thái Bình Dương số hạng
bậc hai đổi dấu, không còn cực đại nội miền, nhất quán với gánh nặng
quốc tế hóa bắt buộc ghi nhận ở nghiên cứu thành phần đảo nhỏ.

Đặc điểm nhà quản trị cấp cao được vận hành hóa từ hai mục khảo sát: số
năm kinh nghiệm của nhà quản trị cấp cao trong ngành (trường
\texttt{b7}) và một biến nhị phân chỉ nhà quản trị cấp cao là nữ (trường
\texttt{b7a} hoặc \texttt{b6a} tùy đợt). Hai biến này được đưa vào mô
hình đa quốc gia để kiểm định hiệu ứng trực tiếp và hiệu ứng điều tiết
của đặc điểm quản trị lên hình dạng đường quan hệ; ở các nghiên cứu
thiếu dữ liệu quản trị, hai biến chỉ được dùng trong phân tích độ vững
chứ không bắt buộc trong mô hình chính, đúng điều kiện sử dụng đã nêu.

\hypertarget{tux1ed5ng-hux1ee3p-biux1ebfn-nghiuxean-cux1ee9u}{%
\subsection{3.3.6 Tổng hợp biến nghiên
cứu}\label{tux1ed5ng-hux1ee3p-biux1ebfn-nghiuxean-cux1ee9u}}

Toàn bộ biến của khung thực nghiệm được tổng hợp trong Bảng 3.1, nêu rõ
với mỗi biến: vai trò trong mô hình, ký hiệu, trường khảo sát WBES
nguồn, cách vận hành hóa và giả thuyết liên quan. Bảng này bảo đảm tính
minh bạch đo lường và khả năng tái lập (Aguinis et al., 2019); chi tiết
vận hành hóa theo từng nghiên cứu thành phần được trình bày trong các
bảng định nghĩa biến ở Mục 3.4.5.

\textbf{Bảng 3.1.} \emph{Tổng hợp biến nghiên cứu của khung thực nghiệm
đa quốc gia.}

\begin{longtable}[]{@{}
  >{\raggedright\arraybackslash}p{(\columnwidth - 8\tabcolsep) * \real{0.2000}}
  >{\raggedright\arraybackslash}p{(\columnwidth - 8\tabcolsep) * \real{0.2000}}
  >{\raggedright\arraybackslash}p{(\columnwidth - 8\tabcolsep) * \real{0.2000}}
  >{\raggedright\arraybackslash}p{(\columnwidth - 8\tabcolsep) * \real{0.2000}}
  >{\raggedright\arraybackslash}p{(\columnwidth - 8\tabcolsep) * \real{0.2000}}@{}}
\toprule\noalign{}
\begin{minipage}[b]{\linewidth}\raggedright
Vai trò
\end{minipage} & \begin{minipage}[b]{\linewidth}\raggedright
Biến (ký hiệu)
\end{minipage} & \begin{minipage}[b]{\linewidth}\raggedright
Trường WBES
\end{minipage} & \begin{minipage}[b]{\linewidth}\raggedright
Vận hành hóa
\end{minipage} & \begin{minipage}[b]{\linewidth}\raggedright
Giả thuyết
\end{minipage} \\
\midrule\noalign{}
\endhead
\bottomrule\noalign{}
\endlastfoot
Biến phụ thuộc & Năng suất lao động, ln(LP) & \texttt{d2}, \texttt{l1} &
Logarit tự nhiên của doanh thu trên số lao động thường xuyên & --- \\
Biến độc lập & Cường độ xuất khẩu (FSTS) & \texttt{d3c} (và
\texttt{d3b}) & Tỷ lệ doanh thu nước ngoài trên tổng doanh thu (0--1),
trung tâm hóa & H1 \\
Biến độc lập & FSTS bình phương (FSTS²) & (suy ra) & Bình phương của
FSTS đã trung tâm hóa, kiểm định phi tuyến & H1, H1b \\
Điều tiết -- tầng doanh nghiệp & Năng lực công nghệ (TCI) & \texttt{b8},
\texttt{e6} & Chỉ số hợp thành chuẩn hóa z (chứng chỉ chất lượng + công
nghệ ngoại nhập) & H2 \\
Điều tiết -- tầng doanh nghiệp & Chấp nhận số (DAI) & \texttt{c22b},
\texttt{k33}/\texttt{k38} & Tầng 1: website (nhị phân); Tầng 2: hợp
thành z (website + thanh toán điện tử) & H3 \\
Điều tiết -- tầng cá nhân & Kinh nghiệm nhà quản trị & \texttt{b7} & Số
năm kinh nghiệm của nhà quản trị cấp cao trong ngành & H4 \\
Điều tiết -- tầng cá nhân & Giới tính nhà quản trị &
\texttt{b7a}/\texttt{b6a} & Biến nhị phân: nhà quản trị cấp cao là nữ &
H4 \\
Điều tiết -- tầng quốc gia & Chế độ thể chế (ICRV) & (phân loại) & Biến
số nguyên 1--6 theo khung ICRV (thể chế giảm dần) & H5 \\
Điều tiết -- thời gian & Đợt khảo sát & (năm khảo sát) & Biến giả
đợt/năm; kiểm định ổn định hệ số Paternoster & H6 \\
Kiểm soát & Quy mô doanh nghiệp & \texttt{l1} & Logarit số lao động
thường xuyên & --- \\
Kiểm soát & Tuổi doanh nghiệp & \texttt{b5} & Năm khảo sát trừ năm thành
lập & --- \\
Kiểm soát & Sở hữu nước ngoài & \texttt{b2b} & Tỷ lệ vốn sở hữu nước
ngoài (hoặc biến nhị phân ≥10\%) & --- \\
Kiểm soát & Hiệu ứng cố định & --- & Ngành (FE), nền kinh tế (FE), năm
(FE) & --- \\
\end{longtable}

\emph{Ghi chú: trường WBES tham chiếu mã biến của bộ công cụ khảo sát;
một số mã thay đổi theo thế hệ lược đồ (PICS3 / Standardized / BREADY)
và được hài hòa hóa theo Phụ lục A. Vận hành hóa chi tiết theo từng
nghiên cứu thành phần xem Mục 3.4.5. Nguồn: tác giả tổng hợp.}

\hypertarget{muxf4-huxecnh-phuxe2n-tuxedch}{%
\section{3.4 Mô hình phân tích}\label{muxf4-huxecnh-phuxe2n-tuxedch}}

\hypertarget{muxf4-huxecnh-phuxe2n-tuxedch-tux1ed5ng-hux1ee3p}{%
\subsection{3.4.1 Mô hình phân tích tổng
hợp}\label{muxf4-huxecnh-phuxe2n-tuxedch-tux1ed5ng-hux1ee3p}}

Mô hình hiệu ứng ngẫu nhiên dựa trên phép biến đổi \(z\) của Fisher:

\[z_i = \frac{1}{2} \ln\left(\frac{1+r_i}{1-r_i}\right)\]

với \(r_i\) là cỡ ảnh hưởng của nghiên cứu \(i\). Cỡ ảnh hưởng gộp được
tính theo trọng số nghịch đảo phương sai (Borenstein et al., 2009).

Dị biệt được đo bằng kiểm định \(Q\) và chỉ số \(I^2\) (Higgins et al.,
2003). Phân tích theo nhóm con và hồi quy phân tích tổng hợp được dùng
để kiểm định các biến điều tiết.

\textbf{Cơ sở kế thừa}: chuẩn phân tích tổng hợp trong quản trị chiến
lược (Combs et al., 2011; Aguinis et al., 2011).

\hypertarget{muxf4-huxecnh-thux1ef1c-nghiux1ec7m-tuyux1ebfn-tuxednh-cux1a1-bux1ea3n}{%
\subsection{3.4.2 Mô hình thực nghiệm: Tuyến tính cơ
bản}\label{muxf4-huxecnh-thux1ef1c-nghiux1ec7m-tuyux1ebfn-tuxednh-cux1a1-bux1ea3n}}

\[\ln(\mathrm{LP})_i = \beta_0 + \beta_1 \mathrm{FSTS}_i + \boldsymbol{\gamma} \mathbf{X}_i + \varepsilon_i\]

với \(\mathbf{X}_i\) là vector control variables.

\hypertarget{muxf4-huxecnh-phi-tuyux1ebfn-h1}{%
\subsection{3.4.3 Mô hình phi tuyến
(H1)}\label{muxf4-huxecnh-phi-tuyux1ebfn-h1}}

\[\ln(\mathrm{LP})_i = \beta_0 + \beta_1 \mathrm{FSTS}_i + \beta_2 \mathrm{FSTS}_i^{2} + \boldsymbol{\gamma} \mathbf{X}_i + \varepsilon_i\]

Nếu \(\beta_1 > 0\) và \(\beta_2 < 0\), chữ U ngược được xác nhận. Điểm
uốn tính theo \(-\beta_1 / (2\beta_2)\).

Đối với Trung Quốc, mô hình giữ nguyên dạng bậc hai trên và được ước
lượng riêng theo từng đợt khảo sát (2012 và 2024) cũng như trên mẫu gộp,
kèm các kiểm định ổn định hệ số xuyên đợt để kiểm tra giả thuyết dị biệt
thời gian H6, thay vì mở rộng lên dạng bậc ba; dạng bậc hai đủ để nhận
dạng điểm uốn và phù hợp với mật độ phân phối cường độ xuất khẩu của mẫu
doanh nghiệp Trung Quốc.

\textbf{Cơ sở kế thừa}: Lu và Beamish (2004), Hitt et al.~(1997),
Contractor et al.~(2003); nghiên cứu thành phần P2 cho dạng bậc hai ở
Trung Quốc với điểm uốn khoảng 47,8\% FSTS, và nghiên cứu thành phần P5
xác nhận dạng bậc hai ổn định theo thời gian với điểm uốn quanh 48\%
FSTS.

\textbf{Kiểm định chữ U của Lind và Mehlum} (Lind \& Mehlum, 2010; Haans
et al., 2016) để xác nhận chặt chẽ chữ U ngược: kiểm định độ dốc dương ở
bên trái và độ dốc âm ở bên phải của miền giá trị FSTS.

\hypertarget{thux1ee7-tux1ee5c-kiux1ec3m-ux111ux1ecbnh-chux1eef-u-ngux1b0ux1ee3c-vuxe0-ux1b0ux1edbc-lux1b0ux1ee3ng-khoux1ea3ng-tin-cux1eady-ux111iux1ec3m-uux1ed1n}{%
\subsection{3.4.3.1 Thủ tục kiểm định chữ U ngược và ước lượng khoảng
tin cậy điểm
uốn}\label{thux1ee7-tux1ee5c-kiux1ec3m-ux111ux1ecbnh-chux1eef-u-ngux1b0ux1ee3c-vuxe0-ux1b0ux1edbc-lux1b0ux1ee3ng-khoux1ea3ng-tin-cux1eady-ux111iux1ec3m-uux1ed1n}}

Sự xuất hiện đồng thời của hệ số bậc nhất dương và hệ số bậc hai âm là
điều kiện cần nhưng chưa đủ để khẳng định quan hệ chữ U ngược, vì cực
trị toán học có thể nằm ngoài miền dữ liệu quan sát, khiến quan hệ thực
tế chỉ đơn điệu trong phạm vi giá trị có ý nghĩa. Luận án vì vậy áp dụng
kiểm định U của Lind và Mehlum (2010) như chuẩn xác nhận chính thức. Thủ
tục đòi hỏi ba điều kiện đồng thời: hệ số bậc nhất dương và hệ số bậc
hai âm; độ dốc của hàm tại biên trái của miền dương và độ dốc tại biên
phải âm; và điểm uốn nằm chặt bên trong khoảng giá trị quan sát của
cường độ xuất khẩu. Kiểm định gộp ba ràng buộc trên thành một kiểm định
giao điểm theo nguyên lý Sasabuchi, cho ra một giá trị p duy nhất; chỉ
khi giá trị p này đủ nhỏ thì hình chữ U ngược mới được xác nhận về mặt
thống kê chứ không chỉ về mặt mô tả.

Điểm uốn được tính theo công thức
\(\text{TP} = -\hat\beta_1 / (2\hat\beta_2)\) trên thang trung tâm hóa
rồi quy đổi về thang gốc bằng cách cộng trung bình mẫu. Vì điểm uốn là
tỷ số phi tuyến của hai hệ số ước lượng, sai số chuẩn của nó không thể
đọc trực tiếp từ bảng hồi quy. Luận án báo cáo khoảng tin cậy 95 phần
trăm theo phương pháp delta (delta method), khai triển tuyến tính bậc
nhất quanh ước lượng điểm để xấp xỉ phương sai của tỷ số (Haans et al.,
2016). Ở các nghiên cứu có mật độ thưa quanh điểm uốn, luận án bổ sung
khoảng tin cậy theo phương pháp lấy mẫu lại bootstrap nhằm nới lỏng giả
định chuẩn của phương pháp delta và phản ánh tính bất đối xứng của phân
phối điểm uốn trong mẫu hữu hạn. Khi kiểm định không bác bỏ tính đơn
điệu trong phạm vi dữ liệu, kết quả được diễn giải là thông tin tích cực
theo khung bão hòa chứ không phải là thất bại của giả thuyết phi tuyến,
vì nó cho biết doanh nghiệp trong mẫu chưa vượt qua ngưỡng nơi lợi suất
quốc tế hóa bắt đầu suy giảm.

Cách vận hành thủ tục này thể hiện rõ ở các nghiên cứu thành phần. Ở
Trung Quốc, kiểm định của Lind và Mehlum được kiểm chứng cho từng đợt
qua điều kiện độ dốc dương có ý nghĩa tại biên trái khi cường độ xuất
khẩu bằng 0 và độ dốc âm có ý nghĩa tại biên phải khi cường độ bằng 1;
điểm uốn được khôi phục bằng phương pháp delta cùng khoảng tin cậy 95
phần trăm theo phổ thể chế của phương pháp delta, cho ra điểm uốn 49,4\%
với khoảng tin cậy từ 43,1\% đến 55,7\% năm 2012 và 47,2\% với khoảng từ
40,8\% đến 53,6\% năm 2024, hai khoảng chồng lấn rộng; trên mẫu gộp điểm
uốn là 48,8\% với khoảng tin cậy từ 44,2\% đến 53,4\%. Ở khung đa quốc
gia, ba điều kiện của Lind và Mehlum, gồm dấu hệ số, điểm uốn nằm chặt
bên trong miền quan sát, và dấu độ dốc tại hai biên, được gộp thành một
kiểm định giao điểm duy nhất. Trường hợp đảo nhỏ Thái Bình Dương minh
họa giá trị thông tin của một kiểm định không bác bỏ: cả hai số hạng mất
ý nghĩa khi đưa vào đồng thời do cộng tuyến cao trong mẫu cường độ thấp
với trung bình cường độ chỉ 0,048, và hệ số bậc hai dương hàm ý dạng chữ
U xuôi thay vì chữ U ngược, một sự đảo dạng bản thân nó nhất quán với
quan hệ đơn điệu âm không có điểm quay nội miền.

\hypertarget{muxf4-huxecnh-ux111iux1ec1u-tiux1ebft-h2h6}{%
\subsection{3.4.4 Mô hình điều tiết
(H2--H6)}\label{muxf4-huxecnh-ux111iux1ec1u-tiux1ebft-h2h6}}

\hypertarget{tux1b0ux1a1ng-tuxe1c-hai-chiux1ec1u-h2-h3}{%
\subsubsection{Tương tác hai chiều (H2,
H3)}\label{tux1b0ux1a1ng-tuxe1c-hai-chiux1ec1u-h2-h3}}

\[\ln(\mathrm{LP})_i = \beta_0 + \beta_1 \mathrm{FSTS}_i + \beta_2 \mathrm{FSTS}_i^{2} + \beta_3 M_i + \beta_4 (\mathrm{FSTS}_i \times M_i) + \beta_5 (\mathrm{FSTS}_i^{2} \times M_i) + \boldsymbol{\gamma} \mathbf{X}_i + \varepsilon_i\]

với \(M_i\) là biến điều tiết (TCI hoặc DAI).

\hypertarget{tux1b0ux1a1ng-tuxe1c-ba-chiux1ec1u-tux1ed5ng-hux1ee3p}{%
\subsubsection{Tương tác ba chiều (tổng
hợp)}\label{tux1b0ux1a1ng-tuxe1c-ba-chiux1ec1u-tux1ed5ng-hux1ee3p}}

\[\ln(\mathrm{LP})_i = \beta_0 + \beta_1 \mathrm{FSTS}_i + \beta_2 \mathrm{FSTS}_i^{2} + \beta_3 D_i + \beta_4 I_i + \beta_5 M_i + \sum \text{(các số hạng tương tác)} + \boldsymbol{\gamma} \mathbf{X}_i + \varepsilon_i\]

với \(D\) = chấp nhận số, \(I\) = thể chế, \(M\) = nhà quản trị.

\textbf{Cơ sở kế thừa}: chuẩn điều tiết trong quản trị chiến lược (Aiken
\& West, 1991; Dawson, 2014).

\textbf{Đóng góp mới}: điều tiết ba chiều (số × thể chế × nhà quản trị)
trong bối cảnh châu Á và Thái Bình Dương với 50 nước chưa được kiểm định
trước đây.

\hypertarget{muxf4-huxecnh-dux1ecb-biux1ec7t-theo-thux1eddi-gian-h6}{%
\subsection{3.4.5 Mô hình dị biệt theo thời gian
(H6)}\label{muxf4-huxecnh-dux1ecb-biux1ec7t-theo-thux1eddi-gian-h6}}

\[\ln(\mathrm{LP})_i = \beta_0 + \beta_1 \mathrm{FSTS}_i + \beta_2 \mathrm{FSTS}_i^{2} + \beta_3 Y_{2024,i} + \beta_4 (\mathrm{FSTS}_i \times Y_{2024,i}) + \beta_5 (\mathrm{FSTS}_i^{2} \times Y_{2024,i}) + \boldsymbol{\gamma} \mathbf{X}_i + \varepsilon_i\]

Nếu \(\beta_4\) hoặc \(\beta_5\) có ý nghĩa, hình dạng đường quan hệ đã
dịch chuyển giữa hai mốc thời gian.

\textbf{Cơ sở kế thừa}: Wu et al.~(2022) với tiến hóa hai mươi năm; Xiao
et al.~(2013) với các điều kiện đặc thù của Trung Quốc; nghiên cứu thành
phần P2 làm cơ sở dạng bậc hai năm 2012.

\textbf{Đóng góp mới}: áp dụng cho Trung Quốc các năm 2012 và 2024 để
kiểm định tính ổn định của điểm uốn dạng bậc hai (khoảng 47,8\% FSTS)
sau hơn một thập kỷ chuyển đổi số.

\hypertarget{kiux1ec3m-ux111ux1ecbnh-ux1ed5n-ux111ux1ecbnh-hux1ec7-sux1ed1-xuyuxean-ux111ux1ee3t-khux1ea3o-suxe1t-theo-paternoster}{%
\subsection{3.4.5.0 Kiểm định ổn định hệ số xuyên đợt khảo sát theo
Paternoster}\label{kiux1ec3m-ux111ux1ecbnh-ux1ed5n-ux111ux1ecbnh-hux1ec7-sux1ed1-xuyuxean-ux111ux1ee3t-khux1ea3o-suxe1t-theo-paternoster}}

Để đánh giá xem hình dạng quan hệ quốc tế hóa--hiệu quả có dịch chuyển
giữa các đợt khảo sát hay không, luận án không chỉ dựa vào biến tương
tác giữa cường độ xuất khẩu và biến giả đợt mà còn áp dụng kiểm định z
của Paternoster và cộng sự (1998) cho tính bằng nhau của hệ số hồi quy
giữa hai mẫu độc lập. Kiểm định này phù hợp với cấu trúc dữ liệu lát cắt
ngang lặp lại, nơi mỗi đợt khảo sát rút một mẫu doanh nghiệp độc lập chứ
không theo dõi cùng doanh nghiệp qua thời gian, khiến so sánh hệ số giữa
các đợt là so sánh giữa hai ước lượng không tương quan. Thống kê kiểm
định nhận dạng chênh lệch hai hệ số chia cho căn bậc hai của tổng phương
sai hai ước lượng, cụ thể
\(z = (\hat\beta_A - \hat\beta_B) / \sqrt{\text{SE}(\hat\beta_A)^2 + \text{SE}(\hat\beta_B)^2}\),
với giá trị p hai phía từ phân phối chuẩn tiêu chuẩn.

Kiểm định được áp dụng riêng cho từng cặp hệ số trọng tâm, gồm số hạng
bậc nhất, số hạng bậc hai và các hệ số trực tiếp của năng lực công nghệ
và chấp nhận số. Việc tách kiểm định theo từng tham số cho phép phân
biệt hai loại dịch chuyển có ý nghĩa lý thuyết khác nhau: dịch chuyển ở
tham số độ cong báo hiệu hình dạng quan hệ thay đổi cấu trúc, còn dịch
chuyển ở hệ số trực tiếp báo hiệu vai trò của một chiều năng lực biến
đổi theo vòng đời số hóa mà không làm đổi dạng đường cong. Khung đa đợt
của luận án vì vậy có thể kết luận, chẳng hạn, rằng ngưỡng tối ưu bền
vững cấu trúc trong khi đóng góp trực tiếp của chấp nhận số lại biến
thiên rõ rệt giữa các đợt. Ở nghiên cứu có sẵn lõi panel nhỏ gồm các
doanh nghiệp xuất hiện ở cả hai đợt, kiểm định Paternoster được bổ trợ
bằng sai số chuẩn phân cụm theo định danh doanh nghiệp nhằm xử lý tương
quan nội nhóm.

Thủ tục này cho ra hai kết cục đối lập về mặt lý thuyết tùy bối cảnh thể
chế. Ở Trung Quốc, kiểm định không bác bỏ bình đẳng hệ số giữa hai đợt
cho cả số hạng bậc nhất và bậc hai, với thống kê dịch chuyển bậc nhất
bằng +0,82 (giá trị p bằng 0,412) và dịch chuyển bậc hai bằng −0,61 (giá
trị p bằng 0,545), hỗ trợ cách diễn giải rằng ngưỡng chữ U ngược là một
quy luật cấu trúc bền vững với điểm uốn ổn định quanh 48\% bất kể thập
niên. Tính bền vững này còn vững trước giới hạn mẫu: trên mẫu con chế
biến chế tạo, điểm uốn là 48,1\% năm 2012 và 46,4\% năm 2024 với các
thống kê Paternoster vẫn không có ý nghĩa. Ngược lại, ở Ấn Độ, kiểm định
Paternoster xuyên đợt bác bỏ bình đẳng hệ số ở mức rất nhỏ hơn 0,0001
cho cả số hạng tuyến tính và bậc hai, một sự sụp đổ trùng với thập niên
thay đổi thể chế ngoại lệ gồm bãi bỏ tiền mặt, thuế hàng hóa và dịch vụ,
luật phá sản, hạ tầng thanh toán hợp nhất, chương trình khuyến khích sản
xuất và đại dịch. Sự đối lập giữa bền vững ở Trung Quốc và sụp đổ ở Ấn
Độ chính là giá trị nhận dạng của kiểm định xuyên đợt: nó phân biệt được
dịch chuyển cấu trúc thực do cú sốc thể chế khỏi biến thiên ngẫu nhiên
của mẫu.

\hypertarget{muxf4-huxecnh-cux1ee5-thux1ec3-nghiuxean-cux1ee9u-3-viux1ec7t-nam-wbes-200920152023}{%
\subsection{3.4.5.1 Mô hình cụ thể: Nghiên cứu 3 (Việt Nam, WBES
2009/2015/2023)}\label{muxf4-huxecnh-cux1ee5-thux1ec3-nghiuxean-cux1ee9u-3-viux1ec7t-nam-wbes-200920152023}}

Nghiên cứu 3 sử dụng chuỗi mô hình lồng nhau M0--M8 ước lượng riêng theo
từng sóng khảo sát (2009, 2015, 2023) và trên mẫu gộp. Ký hiệu tiếng
Việt:

\begin{itemize}
\tightlist
\item
  \textbf{lnNSLD\_it} = lnLP\_it: log năng suất lao động (ln doanh thu
  PPP / lao động thường xuyên)
\item
  \textbf{CDDXK\_c\_it} = FSTS\_c\_it: cường độ xuất khẩu điều chỉnh
  trung bình trong sóng (mean-centred)
\item
  \textbf{CDDXK\_c²\_it} = FSTS\_c²\_it: bình phương cường độ xuất khẩu
  điều chỉnh
\item
  \textbf{NLCN\_z\_it} = TCI\_z\_it: năng lực công nghệ (chuẩn hoá z
  trong sóng)
\item
  \textbf{CSS\_z\_it} = DAI\_z\_it: chỉ số số hoá cơ sở, Tầng 1 (chuẩn
  hoá z trong sóng)
\item
  \textbf{lnLD, TuoiDN, SoHuuNuocNgoai}: kiểm soát quy mô, tuổi, sở hữu
\item
  \textbf{δ\_s}: hiệu ứng cố định ngành (1-digit ISIC); \textbf{λ\_t}:
  hiệu ứng cố định sóng (chỉ pooled)
\end{itemize}

\textbf{M0, Mô hình cơ sở (biến kiểm soát):}
\[\ln\mathrm{NSLD}_{it} = \alpha + \gamma_1 \ln\mathrm{LD}_{it} + \gamma_2 \mathrm{TuoiDN}_{it} + \gamma_3 \mathrm{SoHuuNN}_{it} + \delta_s + [\lambda_t] + \varepsilon_{it}\]

\textbf{M1, Quốc tế hoá tuyến tính:}
\[\ln\mathrm{NSLD}_{it} = \alpha + \beta_1 \mathrm{CDDXK}^{c}_{it} + \boldsymbol{\gamma}\mathbf{X}_{it} + \delta_s + [\lambda_t] + \varepsilon_{it}\]

\textbf{M2, Quốc tế hoá phi tuyến (kiểm định H1, chữ U ngược):}
\[\ln\mathrm{NSLD}_{it} = \alpha + \beta_1 \mathrm{CDDXK}^{c}_{it} + \beta_2 (\mathrm{CDDXK}^{c}_{it})^{2} + \boldsymbol{\gamma}\mathbf{X}_{it} + \delta_s + [\lambda_t] + \varepsilon_{it}\]

H1: \(\beta_1 > 0\) và \(\beta_2 < 0\); điểm quay
\(TP^* = -\beta_1 / (2\beta_2)\), xác nhận bởi kiểm định Lind--Mehlum
(2010).

\textbf{M3, Điều tiết NLCN (kiểm định H2):} \[\begin{aligned}
\ln\mathrm{NSLD}_{it} = {} & \alpha + \beta_1 \mathrm{CDDXK}^{c} + \beta_2 (\mathrm{CDDXK}^{c})^{2} + \beta_3 \mathrm{NLCN}_{z} \\
& + \beta_4(\mathrm{CDDXK}^{c} \times \mathrm{NLCN}_{z}) + \beta_5((\mathrm{CDDXK}^{c})^{2} \times \mathrm{NLCN}_{z}) \\
& + \boldsymbol{\gamma}\mathbf{X} + \delta_s + [\lambda_t] + \varepsilon
\end{aligned}\]

\textbf{M4, Điều tiết CSS/DAI (H3 thăm dò):} \[\begin{aligned}
\ln\mathrm{NSLD}_{it} = {} & \alpha + \beta_1 \mathrm{CDDXK}^{c} + \beta_2 (\mathrm{CDDXK}^{c})^{2} + \beta_3 \mathrm{CSS}_{z} \\
& + \beta_4(\mathrm{CDDXK}^{c} \times \mathrm{CSS}_{z}) + \beta_5((\mathrm{CDDXK}^{c})^{2} \times \mathrm{CSS}_{z}) \\
& + \boldsymbol{\gamma}\mathbf{X} + \delta_s + [\lambda_t] + \varepsilon
\end{aligned}\]

\textbf{M5, NLCN trực tiếp (không tương tác):}
\[\ln\mathrm{NSLD}_{it} = \alpha + \beta_1 \mathrm{CDDXK}^{c} + \beta_2 (\mathrm{CDDXK}^{c})^{2} + \beta_3 \mathrm{NLCN}_{z} + \boldsymbol{\gamma}\mathbf{X} + \delta_s + [\lambda_t] + \varepsilon\]

\textbf{M6, CSS/DAI trực tiếp (không tương tác):}
\[\ln\mathrm{NSLD}_{it} = \alpha + \beta_1 \mathrm{CDDXK}^{c} + \beta_2 (\mathrm{CDDXK}^{c})^{2} + \beta_3 \mathrm{CSS}_{z} + \boldsymbol{\gamma}\mathbf{X} + \delta_s + [\lambda_t] + \varepsilon\]

\textbf{M7, Cả hai trực tiếp, không tương tác:}
\[\ln\mathrm{NSLD}_{it} = \alpha + \beta_1 \mathrm{CDDXK}^{c} + \beta_2 (\mathrm{CDDXK}^{c})^{2} + \beta_3 \mathrm{NLCN}_{z} + \beta_4 \mathrm{CSS}_{z} + \boldsymbol{\gamma}\mathbf{X} + \delta_s + [\lambda_t] + \varepsilon\]

\textbf{M8, Mô hình đầy đủ (trực tiếp + điều tiết CSS):}
\[\begin{aligned}
\ln\mathrm{NSLD}_{it} = {} & \alpha + \beta_1 \mathrm{CDDXK}^{c} + \beta_2 (\mathrm{CDDXK}^{c})^{2} + \beta_3 \mathrm{NLCN}_{z} + \beta_4 \mathrm{CSS}_{z} \\
& + \beta_5(\mathrm{CDDXK}^{c} \times \mathrm{CSS}_{z}) + \beta_6((\mathrm{CDDXK}^{c})^{2} \times \mathrm{CSS}_{z}) \\
& + \boldsymbol{\gamma}\mathbf{X} + \delta_s + [\lambda_t] + \varepsilon
\end{aligned}\]

\emph{Tất cả mô hình đều ước lượng bằng OLS với sai số chuẩn vững HC1.
Hiệu ứng cố định sóng \(\lambda_t\) chỉ áp dụng trong mô hình gộp. Mô
hình được ước lượng riêng cho 3 sóng và gộp (N = 989 / 956 / 1.013 /
2.958).}

\textbf{Bảng định nghĩa biến, Nghiên cứu 3 (Việt Nam):}

\begin{longtable}[]{@{}
  >{\raggedright\arraybackslash}p{(\columnwidth - 6\tabcolsep) * \real{0.2500}}
  >{\raggedright\arraybackslash}p{(\columnwidth - 6\tabcolsep) * \real{0.2500}}
  >{\raggedright\arraybackslash}p{(\columnwidth - 6\tabcolsep) * \real{0.2500}}
  >{\raggedright\arraybackslash}p{(\columnwidth - 6\tabcolsep) * \real{0.2500}}@{}}
\toprule\noalign{}
\begin{minipage}[b]{\linewidth}\raggedright
Ký hiệu
\end{minipage} & \begin{minipage}[b]{\linewidth}\raggedright
Biến WBES
\end{minipage} & \begin{minipage}[b]{\linewidth}\raggedright
Cách tính
\end{minipage} & \begin{minipage}[b]{\linewidth}\raggedright
Vai trò
\end{minipage} \\
\midrule\noalign{}
\endhead
\bottomrule\noalign{}
\endlastfoot
lnNSLD & d2, l1 & ln(d2 / l1): ln(doanh thu PPP / lao động thường xuyên)
& Biến phụ thuộc \\
CDDXK & d3c & d3c / 100: cường độ xuất khẩu trực tiếp (0--1) & Biến độc
lập \\
CDDXK\_c & d3c & CDDXK − trung bình sóng: chuẩn hoá điều chỉnh & Biến
độc lập (centred) \\
CDDXK\_c² & d3c & CDDXK\_c bình phương: hạng phi tuyến & Kiểm định
inverted-U (H1) \\
NLCN\_z & b8, e6 & chuẩn hóa z trong sóng của TB(b8₀₁, e6₀₁): chứng chỉ
chất lượng + công nghệ ngoại & Năng lực công nghệ (H2) \\
CSS\_z & c22b & chuẩn hóa z trong sóng của c22b₀₁: hiện diện website &
Số hoá Tầng 1 (H3 thăm dò) \\
lnLD & l1 & ln(l1): log số lao động thường xuyên & Kiểm soát: quy mô
doanh nghiệp \\
TuoiDN & b5 & năm khảo sát − b5: số năm hoạt động & Kiểm soát: tuổi
doanh nghiệp \\
SoHuuNN & b2b & 1 nếu b2b \textgreater{} 0: có vốn nước ngoài & Kiểm
soát: hình thức sở hữu \\
δ\_s & a4b / a4a & ngành 1-chữ số ISIC (a4b cho 2009/2015; a4a cho 2023)
& Hiệu ứng cố định ngành \\
λ\_t & đợt & chỉ báo sóng khảo sát (chỉ pooled) & Hiệu ứng cố định thời
kỳ \\
\end{longtable}

\hypertarget{muxf4-huxecnh-cux1ee5-thux1ec3-nghiuxean-cux1ee9u-4-singapore-wbes-2023}{%
\subsection{3.4.5.2 Mô hình cụ thể: Nghiên cứu 4 (Singapore, WBES
2023)}\label{muxf4-huxecnh-cux1ee5-thux1ec3-nghiuxean-cux1ee9u-4-singapore-wbes-2023}}

Nghiên cứu 4 sử dụng dữ liệu mặt cắt ngang WBES Singapore 2023 (N =
623). Thiết kế đơn sóng, không có thành phần λ\_t. Ký hiệu tiếng Việt:

\begin{itemize}
\tightlist
\item
  \textbf{lnNSLD\_i} = lnLP\_i: log năng suất lao động (ln doanh thu PPP
  / lao động thường xuyên)
\item
  \textbf{CDDXK\_i} = FSTS\_i: cường độ xuất khẩu trực tiếp (d3c / 100)
\item
  \textbf{CDDXK\_c\_i} = FSTS\_c\_i: CDDXK trung bình mẫu; CDDXK\_c² là
  bình phương
\item
  \textbf{NLCN\_z\_i} = TCI\_z\_i: năng lực công nghệ, chuẩn hóa z của
  trung bình(b8₀₁, e6₀₁)
\item
  \textbf{CSS\_z\_i} = DAI\_z\_i: chỉ số số hoá \textbf{Tầng 1 và 2},
  chuẩn hóa z của trung bình(c22b₀₁, k33₀₁, k38₀₁); khác P3 ở chỗ bao
  gồm thanh toán điện tử hai chiều (Tầng 2)
\item
  \textbf{lnLD\_i}, \textbf{TuoiDN\_i}, \textbf{SoHuuNN\_i}: biến kiểm
  soát tương tự P3
\item
  \textbf{δ\_s}: ngành hiệu ứng cố định (ISIC 1-chữ số)
\item
  \textbf{IMR\_i} = nghịch đảo tỷ lệ Mills từ mô hình probit tham gia
  xuất khẩu (kiểm tra độ nhạy Heckman)
\end{itemize}

\textbf{Chuỗi mô hình lồng nhau M0--M5:}

\begin{quote}
M0 (Baseline): lnNSLD\_i = α + γ₁ lnLD\_i + γ₂ TuoiDN\_i + γ₃ SoHuuNN\_i
+ δ\_s + ε\_i
\end{quote}

\begin{quote}
M1 (Tuyến tính FSTS): lnNSLD\_i = α + β₁ CDDXK\_c\_i + γ·X\_i + δ\_s +
ε\_i
\end{quote}

\begin{quote}
M2 (Bậc hai FSTS, kiểm định H1): lnNSLD\_i = α + β₁ CDDXK\_c\_i + β₂
CDDXK\_c²\_i + γ·X\_i + δ\_s + ε\_i H1: β₁ \textgreater{} 0, β₂
\textless{} 0; TP* = −β₁/(2β₂) ≈ 88,6\% FSTS {[}CI bootstrap {[}53\%,
253\%{]}{]} Lưu ý: Lind--Mehlum p = 0,303, không bác bỏ tuyến tính trong
phạm vi dữ liệu quan sát; đây là kết quả thông tin tích cực theo khung
bão hòa (saturation framework)
\end{quote}

\begin{quote}
M3 (+ NLCN, H1-TCI trực tiếp): lnNSLD\_i = α + β₁ CDDXK\_c + β₂
CDDXK\_c² + β₃ NLCN\_z + γ·X + δ\_s + ε
\end{quote}

\begin{quote}
M4 (+ CSS trực tiếp): lnNSLD\_i = α + β₁ CDDXK\_c + β₂ CDDXK\_c² + β₃
NLCN\_z + β₄ CSS\_z + γ·X + δ\_s + ε
\end{quote}

\begin{quote}
M5 (Mô hình đầy đủ, kiểm định H3): lnNSLD\_i = α + β₁ CDDXK\_c + β₂
CDDXK\_c² + β₃ NLCN\_z + β₄ CSS\_z + β₅(CDDXK\_c × CSS\_z) +
β₆(CDDXK\_c² × CSS\_z) + γ·X + δ\_s + ε H3: β₆ \textgreater{} 0, DAI
khuếch đại lợi nhuận quốc tế hóa và hiệu quả ở cường độ xuất khẩu cao
(coordination platform mechanism; Stallkamp \& Schotter, 2021) Kết quả:
β₆ = +3,119 (p = 0,005) trong mẫu đầy đủ; β₆ = +2,821 (p = 0,003,
F-test) trong mẫu chỉ xuất khẩu (N = 84, lưu ý: công suất thống kê ≈
16\%)
\end{quote}

\textbf{Bảng định nghĩa biến, Nghiên cứu 4 (Singapore):}

\begin{longtable}[]{@{}
  >{\raggedright\arraybackslash}p{(\columnwidth - 6\tabcolsep) * \real{0.2500}}
  >{\raggedright\arraybackslash}p{(\columnwidth - 6\tabcolsep) * \real{0.2500}}
  >{\raggedright\arraybackslash}p{(\columnwidth - 6\tabcolsep) * \real{0.2500}}
  >{\raggedright\arraybackslash}p{(\columnwidth - 6\tabcolsep) * \real{0.2500}}@{}}
\toprule\noalign{}
\begin{minipage}[b]{\linewidth}\raggedright
Ký hiệu VN
\end{minipage} & \begin{minipage}[b]{\linewidth}\raggedright
Mã WBES
\end{minipage} & \begin{minipage}[b]{\linewidth}\raggedright
Cách tính
\end{minipage} & \begin{minipage}[b]{\linewidth}\raggedright
Vai trò
\end{minipage} \\
\midrule\noalign{}
\endhead
\bottomrule\noalign{}
\endlastfoot
lnNSLD & d2, l1 & ln(d2 / l1) & Biến phụ thuộc \\
CDDXK\_c & d3c & FSTS − mean(FSTS): centred & Biến độc lập \\
CDDXK\_c² & d3c & CDDXK\_c bình phương & Phi tuyến (H1) \\
NLCN\_z & b8, e6 & chuẩn hóa z của TB(b8₀₁, e6₀₁) & Năng lực công nghệ
(H2) \\
CSS\_z & c22b, k33, k38 & chuẩn hóa z của TB(c22b₀₁, k33₀₁, k38₀₁):
\textbf{Tầng 1 và 2} & Số hoá Tầng 1 và 2 (H3) \\
lnLD & l1 & ln(l1) & Quy mô doanh nghiệp \\
TuoiDN & b5 & năm khảo sát − b5 & Tuổi doanh nghiệp \\
SoHuuNN & b2b & 1 nếu b2b \textgreater{} 0 & Hình thức sở hữu \\
IMR & probit selection & Nghịch đảo tỷ lệ Mills (kiểm tra Heckman) &
Kiểm soát chọn lựa tham gia xuất khẩu \\
δ\_s & a4b & ngành FE (ISIC 1-chữ số) & Hiệu ứng cố định ngành \\
\end{longtable}

\begin{center}\rule{0.5\linewidth}{0.5pt}\end{center}

\hypertarget{muxf4-huxecnh-cux1ee5-thux1ec3-nghiuxean-cux1ee9u-5-trung-quux1ed1c-wbes-2012-vuxe0-2024}{%
\subsection{3.4.5.3 Mô hình cụ thể: Nghiên cứu 5 (Trung Quốc, WBES 2012
và
2024)}\label{muxf4-huxecnh-cux1ee5-thux1ec3-nghiuxean-cux1ee9u-5-trung-quux1ed1c-wbes-2012-vuxe0-2024}}

Nghiên cứu 5 sử dụng hai sóng WBES Trung Quốc (2012: N = 2.610; 2024: N
= 1.934; pooled N = 4.544). Thiết kế đa sóng không panel (217 doanh
nghiệp xuất hiện cả hai sóng tạo thành ``panel core'' được sử dụng cho
cluster-robust SE). Ký hiệu tiếng Việt:

\begin{itemize}
\tightlist
\item
  \textbf{lnNSLD\_it} = lnLP\_it: log năng suất lao động (biến phụ
  thuộc)
\item
  \textbf{CDDXK\_c\_it} = FSTS\_c\_it: cường độ xuất khẩu centred theo
  trung bình sóng
\item
  \textbf{NLCN\_z\_it} = TCI\_full\_z\_it: năng lực công nghệ toàn diện,
  chuẩn hóa z của TB(b8₀₁, e6₀₁, h1₀₁, h8₀₁); mở rộng hơn P3
\item
  \textbf{CSS\_z\_it} = DAI\_core\_it: chỉ số số hoá cơ bản, chỉ c22b₀₁
  (Tầng 1 mỏng, một chỉ báo nhị phân)
\item
  \textbf{lnLD\_it}, \textbf{TuoiDN\_it}, \textbf{SoHuuNN\_it}: biến
  kiểm soát
\item
  \textbf{δ\_s}: ngành hiệu ứng cố định; \textbf{λ\_t}: đợt hiệu ứng cố
  định (pooled)
\end{itemize}

\textbf{Chuỗi mô hình lồng nhau M0--M6:}

\begin{quote}
M0 (Baseline): lnNSLD\_it = α + γ₁ lnLD\_it + γ₂ TuoiDN\_it + γ₃
SoHuuNN\_it + δ\_s + λ\_t + ε\_it
\end{quote}

\begin{quote}
M1 (Tuyến tính FSTS): lnNSLD\_it = α + β₁ CDDXK\_c\_it + γ·X\_it + δ\_s
+ λ\_t + ε\_it
\end{quote}

\begin{quote}
M2 (Bậc hai FSTS, kiểm định H1): lnNSLD\_it = α + β₁ CDDXK\_c\_it + β₂
CDDXK\_c²\_it + γ·X\_it + δ\_s + λ\_t + ε\_it H1: β₁ \textgreater{} 0,
β₂ \textless{} 0; TP* = −β₁/(2β₂), điểm uốn 49,4\% (2012), 47,2\%
(2024), và 48,8\% (pooled) Kiểm định ổn định cấu trúc: Paternoster et
al.~(1998) z-test: z(FSTS) = +0,82 (p = 0,412); z(FSTS²) = −0,61 (p =
0,545), không bác bỏ bình đẳng hệ số giữa hai sóng
\end{quote}

\begin{quote}
M3 (+ NLCN trực tiếp, H2 level-shift): lnNSLD\_it = α + β₁ CDDXK\_c + β₂
CDDXK\_c² + β₃ NLCN\_z + γ·X + δ\_s + λ\_t + ε Kết quả: β₃ = +0,28
(2012, p \textless{} 0,001), +0,43 (2024, p \textless{} 0,001), TCI là
``bộ tăng mức'' (level-shifter), không điều tiết độ cong
\end{quote}

\begin{quote}
M4 (+ CSS trực tiếp): lnNSLD\_it = α + β₁ CDDXK\_c + β₂ CDDXK\_c² + β₃
NLCN\_z + β₄ CSS\_z + γ·X + δ\_s + λ\_t + ε
\end{quote}

\begin{quote}
M5 (Kiểm định ổn định xuyên sóng, H2b): Ước lượng M2 riêng biệt cho 2012
và 2024, sau đó áp dụng: z = (β̂₁,2012 − β̂₁,2024) / √(SE²₁,2012 +
SE²₁,2024) (Paternoster et al., 1998), tương tự cho β₂ (FSTS²)
\end{quote}

\begin{quote}
M6 (Điều tiết ba chiều, kiểm định H3/H4a/H4b): lnNSLD\_it = α + β₁
CDDXK\_c + β₂ CDDXK\_c² + β₃ NLCN\_z + β₄ CSS\_z + β₅(CDDXK\_c ×
NLCN\_z) + β₆(CDDXK\_c² × NLCN\_z) + β₇(CDDXK\_c × wave) + β₈(CDDXK\_c²
× wave) + γ·X + δ\_s + λ\_t + ε F-tests: F1 (dịch chuyển độ cong xuyên
sóng), F2 (điều tiết năng lực), F3 (dịch chuyển điều kiện) Kết quả: F2 =
3,26 (p = 0,039, không qua Bonferroni α* = 0,017), do đó H4b (không có
điều tiết độ cong theo năng lực) được chấp nhận
\end{quote}

\textbf{Bảng định nghĩa biến, Nghiên cứu 5 (Trung Quốc):}

\begin{longtable}[]{@{}
  >{\raggedright\arraybackslash}p{(\columnwidth - 6\tabcolsep) * \real{0.2500}}
  >{\raggedright\arraybackslash}p{(\columnwidth - 6\tabcolsep) * \real{0.2500}}
  >{\raggedright\arraybackslash}p{(\columnwidth - 6\tabcolsep) * \real{0.2500}}
  >{\raggedright\arraybackslash}p{(\columnwidth - 6\tabcolsep) * \real{0.2500}}@{}}
\toprule\noalign{}
\begin{minipage}[b]{\linewidth}\raggedright
Ký hiệu VN
\end{minipage} & \begin{minipage}[b]{\linewidth}\raggedright
Mã WBES
\end{minipage} & \begin{minipage}[b]{\linewidth}\raggedright
Cách tính
\end{minipage} & \begin{minipage}[b]{\linewidth}\raggedright
Vai trò
\end{minipage} \\
\midrule\noalign{}
\endhead
\bottomrule\noalign{}
\endlastfoot
lnNSLD & d2, l1 & ln(d2 / l1) & Biến phụ thuộc \\
CDDXK\_c & d3c & FSTS − mean\_wave(FSTS) & Biến độc lập \\
CDDXK\_c² & d3c & CDDXK\_c bình phương & Phi tuyến (H1) \\
NLCN\_z & b8, e6, h1, h8 & chuẩn hóa z của TB(b8₀₁, e6₀₁, h1₀₁, h8₀₁):
TCI toàn diện & Năng lực công nghệ (H2) \\
CSS\_z & c22b & c22b₀₁: \textbf{Tầng 1 đơn biến} (binary) & Số hoá Tầng
1 mỏng (kiểm soát) \\
lnLD & l1 & ln(l1) & Quy mô doanh nghiệp \\
TuoiDN & b5 & năm khảo sát − b5 & Tuổi doanh nghiệp \\
SoHuuNN & b2b & 1 nếu b2b \textgreater{} 0 & Hình thức sở hữu \\
δ\_s & ISIC ngành & ngành FE & Hiệu ứng cố định ngành \\
λ\_t & đợt (2012/2024) & đợt FE (chỉ pooled) & Hiệu ứng cố định thời
kỳ \\
\end{longtable}

\emph{Ghi chú: 217 doanh nghiệp xuất hiện cả hai sóng (``panel core'')
tạo thêm nhận dạng biến thiên trong mẫu trong khi clustered SE theo
idstd xử lý tương quan nội nhóm. DAI Trung Quốc chỉ là Tầng 1 đơn biến,
không đủ để kiểm định CDCM; hàm ý chính sách và lý thuyết khác với P4
Singapore (Tầng 1 và 2).}

\begin{center}\rule{0.5\linewidth}{0.5pt}\end{center}

\hypertarget{muxf4-huxecnh-cux1ee5-thux1ec3-nghiuxean-cux1ee9u-7-ux111a-quux1ed1c-gia-chuxe2u-uxe1-wbes-20062026}{%
\subsection{3.4.5.4 Mô hình cụ thể: Nghiên cứu 7 (Đa quốc gia châu Á,
WBES
2006--2026)}\label{muxf4-huxecnh-cux1ee5-thux1ec3-nghiuxean-cux1ee9u-7-ux111a-quux1ed1c-gia-chuxe2u-uxe1-wbes-20062026}}

Nghiên cứu 7 là kiểm định quy mô lớn nhất của luận án, sử dụng dữ liệu
WBES từ \textbf{50 nền kinh tế châu Á và Thái Bình Dương} (103 cặp nền
kinh tế × năm, 2006--2026, gồm Nhật Bản 2025). Mẫu phân tích đạt N =
81.022 (M2) đến N = 79.080 (M5 với kiểm soát đầy đủ). Thiết kế hồi quy
hiệu ứng cố định hai chiều (nền kinh tế và năm) với sai số chuẩn cụm
theo nền kinh tế. Ký hiệu tiếng Việt:

\begin{itemize}
\tightlist
\item
  \textbf{lnNSLD\_it}: ln(năng suất lao động) = ln(doanh thu / lao động
  thường trực), biến phụ thuộc
\item
  \textbf{CDDXK\_c\_it}: cường độ xuất khẩu mean-centred, FSTS\_c =
  (d3c/100) − mean\_pool(FSTS)
\item
  \textbf{CDDXK\_c²\_it}: CDDXK\_c bình phương, kiểm định phi tuyến
\item
  \textbf{NLCN\_z\_it}: năng lực công nghệ z-standardised, TCI\_z =
  chuẩn hóa z(b8, e6)
\item
  \textbf{CSS\_z\_it}: chỉ số số hoá z-standardised, DAI\_z = chuẩn hóa
  z(website + e-pay, Tầng 1 và 2: c22b/e1, k33)
\item
  \textbf{KNQLy\_it}: kinh nghiệm nhà quản lý trong ngành (năm, b7)
\item
  \textbf{NQL\_nu\_it}: top manager là nữ (binary, b7a/b6a)
\item
  \textbf{ICRV\_j}: phân loại chế độ thể chế ICRV (integer 1--6,
  Advanced sang SIDS)
\item
  \textbf{lnLD\_it}: ln(lao động thường trực), kiểm soát quy mô doanh
  nghiệp
\item
  \textbf{TuoiDN\_it}: tuổi doanh nghiệp (năm\_khảo\_sát − b5)
\item
  \textbf{SoHuuNN\_it}: tỷ lệ sở hữu nước ngoài (b6a/100)
\item
  \textbf{δ\_ct}: country × year hiệu ứng cố định (M5)
\end{itemize}

\textbf{Chuỗi mô hình lồng nhau M0--M11:}

\textbf{M0} (mô hình tuyến tính cơ sở):
\[\ln\mathrm{NSLD}_{it} = \alpha + \beta_1 \mathrm{CDDXK}^{c}_{it} + \varepsilon_{it}\]

\textbf{M2} (phi tuyến, kiểm định H1):
\[\ln\mathrm{NSLD}_{it} = \alpha + \beta_1 \mathrm{CDDXK}^{c}_{it} + \beta_2 (\mathrm{CDDXK}^{c}_{it})^{2} + \varepsilon_{it}\]
\[TP^* = -\beta_1 / (2\beta_2); \quad \text{Lind–Mehlum } p < {,}001;\ TP = 51{,}5\%\ (N = 81.022)\]

\textbf{M3} (M2 + kiểm soát cơ bản):
\[\ln\mathrm{NSLD}_{it} = \alpha + \beta_1 \mathrm{CDDXK}^{c} + \beta_2 (\mathrm{CDDXK}^{c})^{2} + \gamma \cdot X_{it} + \varepsilon_{it}\]

\textbf{M5} (M3 + country-year FE, kiểm định vững mạnh nhất):
\[\ln\mathrm{NSLD}_{it} = \alpha + \beta_1 \mathrm{CDDXK}^{c} + \beta_2 (\mathrm{CDDXK}^{c})^{2} + \gamma \cdot X_{it} + \delta_{ct} + \varepsilon_{it}\]
\[TP = 43{,}6\%\ (N = 79.080;\ \text{Lind–Mehlum } p < {,}001)\]

\textbf{M7} (M3 + điều tiết NLCN, kiểm định H2):
\[\ln\mathrm{NSLD}_{it} = \alpha + \beta_1 \mathrm{CDDXK}^{c} + \beta_2 (\mathrm{CDDXK}^{c})^{2} + \beta_3 \mathrm{NLCN}_{z}\]
\[+ \beta_4 (\mathrm{CDDXK}^{c} \times \mathrm{NLCN}_{z}) + \beta_5 ((\mathrm{CDDXK}^{c})^{2} \times \mathrm{NLCN}_{z}) + \gamma \cdot X + \varepsilon\]

\textbf{M8} (M7 + điều tiết CSS, kiểm định H3):
\[\ln\mathrm{NSLD}_{it} = \alpha + \beta_1 \mathrm{CDDXK}^{c} + \beta_2 (\mathrm{CDDXK}^{c})^{2} + \beta_3 \mathrm{NLCN}_{z} + \beta_6 \mathrm{CSS}_{z}\]
\[+ \beta_7 (\mathrm{CDDXK}^{c} \times \mathrm{CSS}_{z}) + \beta_8 ((\mathrm{CDDXK}^{c})^{2} \times \mathrm{CSS}_{z}) + \gamma \cdot X + \varepsilon\]
\[\begin{aligned}
\text{H3: } & \hat\beta_7 = -0{,}271\ (p = {,}149);\ \hat\beta_8 = +0{,}252\ (p = {,}367), \\
& \text{cùng chiều nén nhưng không ý nghĩa toàn mẫu; hiệu ứng mức DAI} = +0{,}201\ (p < {,}001)
\end{aligned}\]

\textbf{M9} (M8 + đặc điểm nhà quản lý, kiểm định H4):
\[+ \beta_9 \mathrm{KNQLy} + \beta_{10} \mathrm{NQL}_{\text{nu}} + \beta_{11}(\mathrm{CDDXK}^{c} \times \mathrm{KNQLy}) + \gamma \cdot X + \varepsilon\]
\[\hat\beta_{10} = -0{,}104^{***} \ (p < {,}001) \text{ (nữ quản lý cấp cao: hệ số mức âm; tương tác với FSTS không ý nghĩa)}\]

\textbf{M10} (M3 + điều tiết ICRV, kiểm định H5):
\[\ln\mathrm{NSLD}_{it} = \alpha + \beta_1 \mathrm{CDDXK}^{c} + \beta_2 (\mathrm{CDDXK}^{c})^{2} + \beta_{11} \mathrm{ICRV}_{j}\]
\[+ \beta_{12}(\mathrm{CDDXK}^{c} \times \mathrm{ICRV}) + \beta_{13}((\mathrm{CDDXK}^{c})^{2} \times \mathrm{ICRV}) + \gamma \cdot X + \varepsilon\]
\[\begin{aligned}
\text{H5 (ba vùng): } & \text{chữ U ngược định hình rõ ở Nhóm IV } (TP = 43{,}0\%,\ p_{LM} < {,}001); \\
& \text{gần tuyến tính ở Nhóm I; mất cấu trúc ở Nhóm V–VI}
\end{aligned}\]

\textbf{M11} (full three-way, kiểm định tương tác tổng hợp ba chiều,
tổng hợp H3 và H5; P7):
\[\ln\mathrm{NSLD}_{it} = \alpha + \beta_1 \mathrm{CDDXK}^{c} + \beta_2 (\mathrm{CDDXK}^{c})^{2} + \beta_3 \mathrm{NLCN}_{z} + \beta_6 \mathrm{CSS}_{z} + \beta_{11} \mathrm{ICRV}\]
\[+ \beta_{7}(\mathrm{CDDXK}^{c} \times \mathrm{CSS}_{z}) + \beta_{8}((\mathrm{CDDXK}^{c})^{2} \times \mathrm{CSS}_{z})\]
\[+ \beta_{12}(\mathrm{CDDXK}^{c} \times \mathrm{ICRV}) + \beta_{13}((\mathrm{CDDXK}^{c})^{2} \times \mathrm{ICRV})\]
\[+ \beta_{14}(\mathrm{CDDXK}^{c} \times \mathrm{CSS}_{z} \times \mathrm{ICRV}) + \beta_9 \mathrm{KNQLy} + \beta_{10} \mathrm{NQL}_{\text{nu}} + \gamma \cdot X + \delta_{ct} + \varepsilon\]
\[\begin{aligned}
& \text{Tương tác chế độ-yếu: FSTS} \times \text{Weak} = -0{,}523\ (p = {,}087); \\
& \text{ba chiều không ý nghĩa; } TP\ (M10) = 45{,}9\%,\ \text{LM } p < {,}001
\end{aligned}\]

\textbf{Bảng định nghĩa biến, Nghiên cứu 7 (đa quốc gia)}

\begin{longtable}[]{@{}
  >{\raggedright\arraybackslash}p{(\columnwidth - 6\tabcolsep) * \real{0.2250}}
  >{\raggedright\arraybackslash}p{(\columnwidth - 6\tabcolsep) * \real{0.2750}}
  >{\raggedright\arraybackslash}p{(\columnwidth - 6\tabcolsep) * \real{0.2750}}
  >{\raggedright\arraybackslash}p{(\columnwidth - 6\tabcolsep) * \real{0.2250}}@{}}
\toprule\noalign{}
\begin{minipage}[b]{\linewidth}\raggedright
Ký hiệu
\end{minipage} & \begin{minipage}[b]{\linewidth}\raggedright
WBES item
\end{minipage} & \begin{minipage}[b]{\linewidth}\raggedright
Cách tính
\end{minipage} & \begin{minipage}[b]{\linewidth}\raggedright
Vai trò
\end{minipage} \\
\midrule\noalign{}
\endhead
\bottomrule\noalign{}
\endlastfoot
lnNSLD & d2/n3, l1 & ln(doanh thu/lao động) & Biến phụ thuộc \\
CDDXK\_c & d3c & (d3c/100) − mean\_pool(FSTS): centred & Biến độc lập
(I) \\
CDDXK\_c² & d3c & CDDXK\_c bình phương & I, phi tuyến \\
NLCN\_z & b8, e6 & chuẩn hóa z(cert chất lượng + CN nước ngoài) & Năng
lực công nghệ \\
CSS\_z & c22b/e1, k33 & chuẩn hóa z(website + thanh toán điện tử): Tầng
1 và 2 & Số hoá (DAI) \\
KNQLy & b7 & năm kinh nghiệm nhà quản lý trong ngành & Quản trị \\
NQL\_nu & b7a/b6a & 1 nếu top manager là nữ & Quản trị \\
ICRV & phân loại & integer 1--6 (tiên tiến đổi mới đến SIDS) & Thể
chế \\
lnLD & l1 & ln(lao động thường trực) & Control (quy mô) \\
TuoiDN & b5 & năm\_khảo\_sát − b5 & Control (tuổi DN) \\
SoHuuNN & b6a & tỷ lệ sở hữu nước ngoài (0--1) & Control (sở hữu) \\
δ\_ct & n/a & country × year FE & hiệu ứng cố định \\
\end{longtable}

\emph{Ghi chú: Mẫu phân tích giảm nhẹ từ N = 81.022 (M2) xuống 79.080
(M5, thêm kiểm soát); các mô hình điều tiết M7, M8 và M10 giữ N tương
đương 79.080--79.220. DAI P7 là Tầng 1 và 2 khi k33 có mặt, Tầng 1 only
khi thiếu, khác với P3 Vietnam (chỉ Tầng 1) và đồng nhất với P4
Singapore (Tầng 1 và 2). ICRV là biến điều tiết liên tục (integer),
không phải dummies, nhằm giữ N đủ lớn trong M10--M11.}

\begin{center}\rule{0.5\linewidth}{0.5pt}\end{center}

\hypertarget{muxf4-huxecnh-cux1ee5-thux1ec3-nghiuxean-cux1ee9u-8-ux111ux1ea3o-nhux1ecf-thuxe1i-buxecnh-dux1b0ux1a1ng-n-1.450-nhuxf3m-vi}{%
\subsubsection{3.4.5.5 Mô hình cụ thể: Nghiên cứu 8: đảo nhỏ Thái Bình
Dương (N = 1.450, Nhóm
VI)}\label{muxf4-huxecnh-cux1ee5-thux1ec3-nghiuxean-cux1ee9u-8-ux111ux1ea3o-nhux1ecf-thuxe1i-buxecnh-dux1b0ux1a1ng-n-1.450-nhuxf3m-vi}}

Nghiên cứu 8 phân tích mẫu gồm \textbf{N = 1.450 doanh nghiệp} tại
\textbf{7 nền kinh tế Pacific SIDS} (Fiji, Kiribati, Papua New Guinea,
Samoa, Solomon Islands, Tonga, Vanuatu) từ các sóng WBES 2009--2025. Đây
là nhóm thể chế ICRV Nhóm VI, cấp thấp nhất trong phân loại 6 chế độ,
đặc trưng bởi thị trường nội địa cực nhỏ, chi phí thương mại cao và hỗ
trợ thể chế yếu. Trong tổng mẫu, chỉ \textbf{212 doanh nghiệp có hoạt
động xuất khẩu} (14,6\%), phản ánh cấu trúc thị trường đảo nhỏ.

Phát hiện trọng tâm: \textbf{gánh nặng quốc tế hóa bắt buộc (Forced
Internationalization Penalty, FIP)}. Tại Nhóm VI, dạng hình chữ U ngược
\textbf{mất cấu trúc}: trên mẫu đầy đủ bảy nền (N = 1.450), cả số hạng
tuyến tính lẫn bậc hai đều mất ý nghĩa khi đưa vào đồng thời và cả bốn
điều kiện cần của một đường cong chữ U ngược đều thất bại (xem Chương 4,
Mục 4.7.2 và 4.7.6), nên quan hệ là \textbf{đơn điệu âm về xu hướng},
không có điểm quay nội miền, trái ngược với hình chữ U ngược ở P3--P7.
Hệ số FIP âm \textbf{mạnh, có ý nghĩa} (β₁ = −1,339, p \textless{} ,001)
ghi nhận trên \textbf{phiên bản dữ liệu gốc / bản dựng hạn chế} (N =
959; lõi ba nền có đủ kiểm soát đợt trước 2018: N = 209) và được trình
bày như bằng chứng độ vững nội bộ ở Chương 4 (Mục 4.7.7--4.7.8).

\textbf{Ký hiệu biến (tiếng Việt ↔ mã WBES):}

\begin{longtable}[]{@{}
  >{\raggedright\arraybackslash}p{(\columnwidth - 6\tabcolsep) * \real{0.2000}}
  >{\raggedright\arraybackslash}p{(\columnwidth - 6\tabcolsep) * \real{0.3556}}
  >{\raggedright\arraybackslash}p{(\columnwidth - 6\tabcolsep) * \real{0.2000}}
  >{\raggedright\arraybackslash}p{(\columnwidth - 6\tabcolsep) * \real{0.2444}}@{}}
\toprule\noalign{}
\begin{minipage}[b]{\linewidth}\raggedright
Ký hiệu
\end{minipage} & \begin{minipage}[b]{\linewidth}\raggedright
Tên tiếng Việt
\end{minipage} & \begin{minipage}[b]{\linewidth}\raggedright
Mã WBES
\end{minipage} & \begin{minipage}[b]{\linewidth}\raggedright
Cách tính
\end{minipage} \\
\midrule\noalign{}
\endhead
\bottomrule\noalign{}
\endlastfoot
lnNSLD & Log năng suất lao động & d2, l1 & ln(doanh thu PPP / lao động
thường trực) \\
CDDXK\_c & Cường độ xuất khẩu (điều chỉnh) & d3c & (d3c/100) −
mean\_wave(FSTS): centred theo sóng \\
CDDXK\_c² & CDDXK\_c bình phương & d3c & CDDXK\_c × CDDXK\_c \\
NLCN\_z & Năng lực công nghệ & b8, e6 & chuẩn hóa z(trung bình cert chất
lượng + CN nước ngoài) \\
CSS\_z & Chỉ số số hoá (Tầng 1) & c22b & chuẩn hóa z(website binary):
\textbf{chỉ Tầng 1, không dynamic} \\
lnLD & Log quy mô lao động & l1 & ln(lao động thường trực) \\
TuoiDN & Tuổi doanh nghiệp & b5 & năm\_khảo\_sát − b5 \\
SoHuuNN & Tỷ lệ sở hữu nước ngoài & b6a & tỷ lệ 0--1 \\
δ\_c & Country hiệu ứng cố định & a3b & biến giả theo quốc gia \\
λ\_t & đợt hiệu ứng cố định & đợt & biến giả theo sóng điều tra \\
\end{longtable}

\emph{Ghi chú: CSS\_z (DAI) là biến đại diện tĩnh Tầng 1 (website binary
c22b). WBES không thu thập đủ thang đo cho số hoá động tại các đảo nhỏ,
KHÔNG gọi là ``năng lực số hoá động''.}

\textbf{Chuỗi mô hình M0--M3 (P8):}

M0, cơ sở kiểm soát:

\[\ln\mathrm{NSLD}_{it} = \alpha + \gamma_1 \ln\mathrm{LD}_{it} + \gamma_2 \mathrm{TuoiDN}_{it} + \gamma_3 \mathrm{SoHuuNN}_{it} + \delta_c + \lambda_t + \varepsilon_{it}\]

M1, Kiểm định FIP (H1b): thêm cường độ xuất khẩu tuyến tính

\[\ln\mathrm{NSLD}_{it} = \alpha + \beta_1 \mathrm{CDDXK}^{c}_{it} + \gamma \cdot X_{it} + \delta_c + \lambda_t + \varepsilon_{it}\]

\[\beta_1 = -1{,}339,\ p < {,}001\ (\text{country+year FE, phiên bản dữ liệu gốc, N} = 959) \Rightarrow \textbf{FIP (bản dựng gốc)}\]

M2, Kiểm định phi tuyến: thêm CDDXK\_c²

\[\ln\mathrm{NSLD}_{it} = \alpha + \beta_1 \mathrm{CDDXK}^{c}_{it} + \beta_2 (\mathrm{CDDXK}^{c}_{it})^{2} + \gamma \cdot X_{it} + \delta_c + \lambda_t + \varepsilon_{it}\]

\[\begin{aligned}
& \beta_1\ \text{NS},\ \beta_2\ \text{NS};\ \text{Lind–Mehlum U-test NS} \\
& \Rightarrow \text{không có điểm quay, đơn điệu âm}
\end{aligned}\]

M3, Kiểm định năng lực điều tiết (H2 null cho SIDS):

\[\begin{aligned}
\ln\mathrm{NSLD}_{it} = {} & \alpha + \beta_1 \mathrm{CDDXK}^{c}_{it} + \beta_2 (\mathrm{CDDXK}^{c}_{it})^{2} + \beta_3 \mathrm{NLCN}_{z,it} \\
& + \beta_4 \mathrm{CSS}_{z,it} + \gamma \cdot X_{it} + \delta_c + \lambda_t + \varepsilon_{it}
\end{aligned}\]

\[\begin{aligned}
& \beta_3\ (\mathrm{NLCN}_{z}):\ p = {,}003\ (\text{dương, có ý nghĩa});\ \beta_4\ (\mathrm{CSS}_{z}):\ p = {,}285\ \text{NS} \\
& \Rightarrow \text{TCI nâng mặt bằng năng suất; không có biến điều tiết uốn đường cong (H2 null)}
\end{aligned}\]

\textbf{Phân tích độ vững (kiểm tra độ vững):}

\begin{longtable}[]{@{}llll@{}}
\toprule\noalign{}
Mô hình & β(CDDXK\_c) & SE & p \\
\midrule\noalign{}
\endhead
\bottomrule\noalign{}
\endlastfoot
M1 country+year FE (N = 959) & −1,339 & 0,386 & \textless,001 \\
Year FE only (không country FE) & −3,351 & 0,808 & \textless,001 \\
Bivariate (không controls) & −0,864 & 0,441 & ,050 \\
chỉ doanh nghiệp xuất khẩu (N = 26) & −1,176 & 0,763 & ,130 (NS) \\
\end{longtable}

\emph{Bảng độ vững trên thuộc \textbf{phiên bản dữ liệu gốc} (N = 959);
hệ số âm nhất quán về dấu qua các mô hình, suy luận FIP dựa chủ yếu vào
các mô hình hiệu ứng cố định (M1, year-FE), trong khi chỉ doanh nghiệp
xuất khẩu (N = 26) thiếu công suất thống kê. Trên mẫu đầy đủ bảy nền (N
= 1.450), dạng đơn điệu âm vẫn là xu hướng nhưng các số hạng mất ý nghĩa
thống kê khi đưa vào đồng thời (Chương 4, Mục 4.7.2).}

\textbf{Bảng định nghĩa biến, Nghiên cứu 8 (SIDS)}

\begin{longtable}[]{@{}
  >{\raggedright\arraybackslash}p{(\columnwidth - 6\tabcolsep) * \real{0.2250}}
  >{\raggedright\arraybackslash}p{(\columnwidth - 6\tabcolsep) * \real{0.2750}}
  >{\raggedright\arraybackslash}p{(\columnwidth - 6\tabcolsep) * \real{0.2750}}
  >{\raggedright\arraybackslash}p{(\columnwidth - 6\tabcolsep) * \real{0.2250}}@{}}
\toprule\noalign{}
\begin{minipage}[b]{\linewidth}\raggedright
Ký hiệu
\end{minipage} & \begin{minipage}[b]{\linewidth}\raggedright
WBES item
\end{minipage} & \begin{minipage}[b]{\linewidth}\raggedright
Cách tính
\end{minipage} & \begin{minipage}[b]{\linewidth}\raggedright
Vai trò
\end{minipage} \\
\midrule\noalign{}
\endhead
\bottomrule\noalign{}
\endlastfoot
lnNSLD & d2, l1 & ln(doanh thu PPP / lao động) & Biến phụ thuộc \\
CDDXK\_c & d3c & (d3c/100) − mean\_wave: centred & Biến độc lập (I) \\
CDDXK\_c² & d3c & CDDXK\_c bình phương & I, kiểm định phi tuyến \\
NLCN\_z & b8, e6 & chuẩn hóa z(cert + CN nước ngoài) & Năng lực công
nghệ \\
CSS\_z & c22b & chuẩn hóa z(website binary): Tầng 1 only & Số hoá, biến
đại diện tĩnh \\
lnLD & l1 & ln(lao động thường trực) & Control (quy mô) \\
TuoiDN & b5 & năm\_KS − năm\_thành\_lập & Control (tuổi DN) \\
SoHuuNN & b6a & tỷ lệ sở hữu nước ngoài & Control (sở hữu) \\
δ\_c & a3b & country FE & hiệu ứng cố định \\
λ\_t & đợt & đợt FE & hiệu ứng cố định \\
\end{longtable}

\emph{Ghi chú: FIP là điều kiện biên (boundary condition) của mô hình
chữ U ngược, áp dụng đặc thù tại ICRV Nhóm VI nơi 3 điều kiện cấu trúc
bị vi phạm đồng thời: (1) không có thị trường nội địa đủ lớn, (2) chi
phí thương mại không kiểm soát được, (3) hỗ trợ thể chế không đủ chức
năng. DAI (CSS\_z) Tầng 1 only, WBES không thu thập thang đo đầy đủ cho
số hoá động tại SIDS.}

\begin{center}\rule{0.5\linewidth}{0.5pt}\end{center}

\hypertarget{muxf4-huxecnh-cux1ee5-thux1ec3-nghiuxean-cux1ee9u-9-ux1ea5n-ux111ux1ed9-ux111ux1ed9-bux1ec1n-ngux1b0ux1ee1ng-qua-chuyux1ec3n-ux111ux1ed5i-thux1ec3-chux1ebf-p9}{%
\subsubsection{3.4.5.6 Mô hình cụ thể, Nghiên cứu 9: Ấn Độ, độ bền
ngưỡng qua chuyển đổi thể chế
(P9)}\label{muxf4-huxecnh-cux1ee5-thux1ec3-nghiuxean-cux1ee9u-9-ux1ea5n-ux111ux1ed9-ux111ux1ed9-bux1ec1n-ngux1b0ux1ee1ng-qua-chuyux1ec3n-ux111ux1ed5i-thux1ec3-chux1ebf-p9}}

Nghiên cứu 9 sử dụng ba đợt Khảo sát Doanh nghiệp Ngân hàng Thế giới cho
Ấn Độ (2014 PICS3, \(N_{\text{phân tích}} = 8.941\); 2022 BEE,
\(N_{\text{phân tích}} = 9.300\); 2025 BREADY,
\(N_{\text{phân tích}} = 10.476\); mẫu gộp \(N = 28.717\)), mẫu gộp đơn
quốc gia lớn nhất của luận án. Ấn Độ thuộc Nhóm IV ICRV (Chuyển đổi thu
nhập trung bình--thấp); cửa sổ 2014--2025 bao trùm chuỗi cú sốc thể chế
(bãi bỏ tiền mặt 2016, GST 2017, IBC 2016, UPI 2016, PLI và Atmanirbhar
Bharat 2020, COVID-19), tạo điều kiện stress-test cho giả thuyết bền
vững ngưỡng.

Mô hình ước lượng OLS với sai số chuẩn vững HC1, gồm số hạng FSTS bậc
nhất và bậc hai (centred theo đợt) cùng các tương tác năng lực (TCI,
DAI). Ba kiểm định bổ sung: (i) Lind--Mehlum U-test xác nhận U ngược
theo từng đợt; (ii) z-test Paternoster (1998) so sánh hệ số giữa 2014 và
2025 (cửa sổ 11 năm chính), báo cáo cả sai số chuẩn HC1 và cụm theo
bang; (iii) mô hình gộp ba đợt với biến giả/xu thế đợt khảo sát làm kiểm
định độ vững. Việc UPI ra mắt tháng 4/2016 cung cấp thí nghiệm gần tự
nhiên cho bão hòa DAI; riêng đợt 2025 (BREADY) bổ sung chỉ báo Tầng 2
(\texttt{k33}, tỷ trọng doanh thu qua thanh toán điện tử) cho phép kiểm
định điều tiết DAI Tầng 2, được xử lý như phân tích thăm dò vì chỉ có ở
một đợt.

\hypertarget{sai-sux1ed1-chuux1ea9n}{%
\subsection{3.4.6 Sai số chuẩn}\label{sai-sux1ed1-chuux1ea9n}}

Tất cả các mô hình đều sử dụng HC1/HC3 sai số chuẩn vững (Long \& Ervin,
2000; White, 1980) để xử lý phương sai thay đổi. Đối với multi-country,
bổ sung clustered SE theo country.

\hypertarget{lux1ef1a-chux1ecdn-cux1ea5u-truxfac-sai-sux1ed1-chuux1ea9n-theo-thiux1ebft-kux1ebf-dux1eef-liux1ec7u}{%
\subsection{3.4.6.1 Lựa chọn cấu trúc sai số chuẩn theo thiết kế dữ
liệu}\label{lux1ef1a-chux1ecdn-cux1ea5u-truxfac-sai-sux1ed1-chuux1ea9n-theo-thiux1ebft-kux1ebf-dux1eef-liux1ec7u}}

Cấu trúc sai số chuẩn được lựa chọn theo bản chất thiết kế của từng
nghiên cứu chứ không áp dụng một quy ước đồng nhất. Các nghiên cứu đơn
quốc gia trên lát cắt ngang dùng sai số chuẩn vững với phương sai thay
đổi loại HC1, một hiệu chỉnh chuẩn xử lý phương sai sai số không đồng
nhất giữa các doanh nghiệp mà không đòi hỏi giả định về dạng cụ thể của
phương sai. Khung đa quốc gia bổ sung phân cụm sai số chuẩn theo nền
kinh tế, vì thiết kế chọn mẫu phân tầng của WBES tạo ra tương quan nội
cụm giữa các doanh nghiệp trong cùng một nền kinh tế, khiến giả định độc
lập của HC1 không còn thỏa đáng. Phân cụm theo nền kinh tế cũng phản ánh
đúng cấp độ mà các cú sốc thể chế và chính sách tác động đồng thời lên
nhiều doanh nghiệp.

Cụ thể, mọi mô hình trong khung đa quốc gia đều mang hiệu ứng cố định
hai chiều theo nền kinh tế và theo năm, cùng sai số chuẩn vững phân cụm
theo nền kinh tế theo hiệu chỉnh CRV1. Cấp phân cụm theo nền kinh tế
được chọn vì thiết kế chọn mẫu phân tầng của bộ công cụ tạo tương quan
nội cụm giữa các doanh nghiệp trong cùng một nền, và vì các cú sốc thể
chế cùng chính sách tác động đồng thời lên nhiều doanh nghiệp ở cùng cấp
này.

Một giới hạn cần ghi nhận của phân cụm theo nền kinh tế là số cụm hữu
hạn có thể làm sai số chuẩn tiệm cận kém chính xác khi quy mô cụm chênh
lệch lớn, đặc biệt giữa các nền kinh tế đông doanh nghiệp và các đảo nhỏ
chỉ có vài trăm quan sát. Trong các tình huống này, luận án diễn giải
suy luận một cách thận trọng và neo kết luận chính vào tính nhất quán về
dấu qua nhiều mô hình thay vì vào ý nghĩa của một hệ số đơn lẻ. Việc
luận án không áp dụng trọng số chọn mẫu trong ước lượng chính tuân theo
khuyến nghị của Solon, Haider và Wooldridge (2015): khi mô hình đã kiểm
soát các chiều phân tầng qua biến kiểm soát và qua hiệu ứng cố định, ước
lượng bình phương nhỏ nhất không trọng số cho hệ số nhất quán và thường
hiệu quả hơn, trong khi trọng số được dành cho phân tích độ vững nhằm
kiểm tra tính ổn định.

\hypertarget{kiux1ec3m-ux111ux1ecbnh-ux111ux1ed9-vux1eefng}{%
\section{3.5 Kiểm định độ
vững}\label{kiux1ec3m-ux111ux1ecbnh-ux111ux1ed9-vux1eefng}}

\hypertarget{thay-ux111ux1ed5i-thux1b0ux1edbc-ux111o}{%
\subsection{3.5.1 Thay đổi thước
đo}\label{thay-ux111ux1ed5i-thux1b0ux1edbc-ux111o}}

Lớp độ vững thứ nhất hoán đổi thước đo các biến trọng tâm. Biến phụ
thuộc năng suất lao động dạng logarit được thay bằng tỷ suất lợi nhuận
trên doanh thu, tăng trưởng doanh thu và tăng trưởng việc làm để kiểm
tra liệu hình dạng quan hệ có phụ thuộc vào một định nghĩa hiệu quả duy
nhất hay không. Biến độc lập cường độ xuất khẩu được thay bằng trạng
thái xuất khẩu nhị phân nhằm phân biệt hiệu ứng ở biên tham gia khỏi
hiệu ứng ở biên cường độ. Biến điều tiết số được hoán đổi giữa biến thể
chỉ gồm Tầng 1 và biến thể gồm cả Tầng 2 trên mẫu con từ năm 2024 trở
đi, nơi các chỉ báo số phong phú hơn đã sẵn có.

\hypertarget{mux1eabu-phux1ee5}{%
\subsection{3.5.2 Mẫu phụ}\label{mux1eabu-phux1ee5}}

Lớp thứ hai giới hạn mẫu để loại trừ các giải thích thay thế dựa trên
thành phần mẫu. Các mẫu con gồm: loại doanh nghiệp siêu nhỏ dưới năm lao
động, chỉ giữ doanh nghiệp vừa và nhỏ, chỉ giữ ngành chế biến chế tạo,
chỉ giữ doanh nghiệp có xuất khẩu với cường độ dương, và tách theo từng
chế độ trong sáu chế độ thể chế ICRV. Việc điểm uốn và dấu của độ cong
được bảo toàn qua các mẫu con này là bằng chứng rằng quan hệ trọng tâm
không phải tạo tác của một nhóm doanh nghiệp đặc thù.

\hypertarget{giuxe1-trux1ecb-ngoux1ea1i-lai}{%
\subsection{3.5.3 Giá trị ngoại
lai}\label{giuxe1-trux1ecb-ngoux1ea1i-lai}}

Lớp thứ ba kiểm soát ảnh hưởng của giá trị ngoại lai bằng cách cắt đuôi
năng suất lao động và số lao động tại các ngưỡng phân vị 1 và 99, rồi
lặp lại tại ngưỡng 5 và 95, áp dụng riêng trong từng cặp nền kinh tế và
năm để tránh nhầm chênh lệch mặt bằng giữa các nền kinh tế với giá trị
ngoại lai thực.

\hypertarget{ux111a-cux1ed9ng-tuyux1ebfn-vuxe0-phux1b0ux1a1ng-sai-thay-ux111ux1ed5i}{%
\subsection{3.5.4 Đa cộng tuyến và phương sai thay
đổi}\label{ux111a-cux1ed9ng-tuyux1ebfn-vuxe0-phux1b0ux1a1ng-sai-thay-ux111ux1ed5i}}

Hai mối lo của hồi quy có số hạng phi tuyến và tương tác là đa cộng
tuyến và phương sai sai số thay đổi. Để xử lý đa cộng tuyến giữa số hạng
bậc nhất và số hạng bậc hai của cường độ xuất khẩu, cũng như giữa các
thành phần của số hạng tương tác, biến cường độ xuất khẩu được căn giữa
quanh trung bình mẫu trước khi lập bình phương và lập tương tác; căn
giữa làm giảm mạnh tương quan giả tạo giữa một biến và bình phương của
chính nó mà không làm thay đổi hệ số bậc hai hay vị trí điểm uốn. Hệ số
phóng đại phương sai được tính cho mọi mô hình với ngưỡng cảnh báo dưới
năm (Hair et al., 2019); trong các mô hình đã căn giữa, hệ số phóng đại
phương sai nằm an toàn dưới ngưỡng này, xác nhận rằng các ước lượng hệ
số không bị bất ổn do cộng tuyến. Đối với phương sai sai số thay đổi,
kiểm định Breusch và Pagan (1979) được dùng để phát hiện, và quan trọng
hơn, mọi suy luận đều dựa trên sai số chuẩn vững với phương sai thay đổi
và được gom cụm theo nền kinh tế, do đó kết luận thống kê có hiệu lực
ngay cả khi phương sai sai số không đồng nhất, một cấu trúc sai số phù
hợp với thiết kế dữ liệu gộp nhiều nền kinh tế đã trình bày ở Mục 3.4.6.

\hypertarget{phuxe2n-tuxedch-ux111ux1ed9-nhux1ea1y-cho-phuxe2n-tuxedch-tux1ed5ng-hux1ee3p}{%
\subsection{3.5.5 Phân tích độ nhạy cho phân tích tổng
hợp}\label{phuxe2n-tuxedch-ux111ux1ed9-nhux1ea1y-cho-phuxe2n-tuxedch-tux1ed5ng-hux1ee3p}}

Khung phân tích tổng hợp (P6) áp dụng một bộ kiểm định độ nhạy riêng phù
hợp với cấu trúc dữ liệu ba tầng. Kiểm định bỏ một nghiên cứu lần lượt
loại từng nghiên cứu khỏi mẫu và tái ước lượng cỡ ảnh hưởng gộp, nhằm
bảo đảm kết luận không bị chi phối bởi một nghiên cứu đơn lẻ có trọng số
lớn hay giá trị cực đoan. Lệch lạc công bố được đánh giá bằng nhiều công
cụ bổ trợ: kiểm tra trực quan biểu đồ phễu, kiểm định bất đối xứng của
Egger, và thủ tục cắt và điền của Duval và Tweedie (2000) để ước lượng
và hiệu chỉnh số nghiên cứu khuyết có thể có. Ngoài ra, cấu trúc ba tầng
cho phép so sánh mô hình ba tầng với mô hình hai tầng truyền thống nhằm
kiểm tra xem việc tách riêng phương sai trong nội bộ nghiên cứu khỏi
phương sai giữa các nghiên cứu có thực sự cải thiện độ phù hợp hay
không, đồng thời cho phép phân rã \(I^2\) thành thành phần trong và giữa
nghiên cứu để định vị nguồn dị biệt. Kết hợp các kiểm định này bảo đảm
rằng cỡ ảnh hưởng gộp và kết luận điều tiết theo ICRV là vững trước cả
lựa chọn mẫu nghiên cứu lẫn giả định cấu trúc phương sai.

\hypertarget{giao-thux1ee9c-ux111ux1ed9-vux1eefng-ux111a-lux1edbp-vuxe0-ux111a-ux111ux1ee3t-cho-khung-thux1ef1c-nghiux1ec7m}{%
\subsection{3.5.5.1 Giao thức độ vững đa lớp và đa đợt cho khung thực
nghiệm}\label{giao-thux1ee9c-ux111ux1ed9-vux1eefng-ux111a-lux1edbp-vuxe0-ux111a-ux111ux1ee3t-cho-khung-thux1ef1c-nghiux1ec7m}}

Khung thực nghiệm áp dụng một giao thức độ vững có cấu trúc nhằm bảo đảm
suy luận trọng tâm không phụ thuộc vào một quyết định đo lường, một mẫu
con, hay một mô hình đơn lẻ. Lớp thứ nhất hoán đổi thước đo: biến phụ
thuộc được thay bằng tỷ suất lợi nhuận trên doanh thu, tăng trưởng doanh
thu và tăng trưởng việc làm; biến độc lập được thay bằng trạng thái xuất
khẩu nhị phân; biến điều tiết số được hoán đổi giữa biến thể Tầng 1 và
biến thể Tầng 1 cộng Tầng 2 trên mẫu con có đủ chỉ báo. Lớp thứ hai giới
hạn mẫu con: loại doanh nghiệp siêu nhỏ, chỉ giữ doanh nghiệp vừa và
nhỏ, chỉ giữ ngành chế tạo, chỉ giữ doanh nghiệp có xuất khẩu, và tách
theo từng chế độ thể chế. Lớp thứ ba kiểm soát phân phối qua cắt đuôi
tại hai ngưỡng phân vị khác nhau trong từng cặp nền kinh tế và năm.

Lớp thứ tư kiểm tra các giả định kinh tế lượng phụ trợ. Hệ số phóng đại
phương sai (variance inflation factor, VIF) được tính cho mọi mô hình
với ngưỡng cảnh báo dưới năm để phát hiện đa cộng tuyến đe dọa độ ổn
định hệ số; kiểm định Breusch--Pagan đánh giá phương sai thay đổi. Đối
với cấu phần phân tích tổng hợp, giao thức bổ sung kiểm định loại từng
nghiên cứu (leave-one-out) tái tính hiệu ứng gộp nhằm phát hiện ảnh
hưởng quá mức của một nghiên cứu đơn lẻ, cùng thủ tục cắt và điền
(trim-and-fill) đánh giá thiên lệch công bố. Luận án không áp dụng hiệu
chỉnh kiểm định bội chính thức cho toàn bộ bảng độ vững, vì các bảng
kiểm tra các quan ngại nhận diện khác nhau thay vì lặp lại cùng một giả
thuyết; thay vào đó, suy luận thực chất được neo vào mẫu hình nhất quán
qua các bảng và tính nhất quán về dấu của ước lượng trọng tâm chứ không
vào ý nghĩa của một bảng độ vững đơn lẻ.

\hypertarget{chiux1ebfn-lux1b0ux1ee3c-nhux1eadn-dux1ea1ng-vuxe0-ranh-giux1edbi-suy-luux1eadn-nhuxe2n-quux1ea3}{%
\subsection{3.5.6 Chiến lược nhận dạng và ranh giới suy luận nhân
quả}\label{chiux1ebfn-lux1b0ux1ee3c-nhux1eadn-dux1ea1ng-vuxe0-ranh-giux1edbi-suy-luux1eadn-nhuxe2n-quux1ea3}}

Vì dữ liệu chủ đạo của luận án là lát cắt ngang lặp lại (repeated
cross-sections) chứ không phải bảng cân bằng theo dõi cùng doanh nghiệp
qua thời gian, luận án thừa nhận minh bạch các ranh giới suy luận và áp
dụng một chiến lược nhận dạng nhiều lớp để củng cố tính giá trị nội tại
trong giới hạn cho phép. Thứ nhất, mọi mô hình đa quốc gia đều bao gồm
hiệu ứng cố định hai chiều (nền kinh tế và năm), qua đó hấp thụ toàn bộ
thành phần bất biến theo nước (tiền tệ, mặt bằng giá, chế độ kế toán,
đặc điểm thể chế ổn định) và các cú sốc chung theo năm; hệ số quan hệ
quốc tế hóa và hiệu quả do đó được nhận dạng từ biến thiên nội bộ nền
kinh tế--năm, không bị nhiễu bởi khác biệt cấu trúc cố định giữa các
nước (Wooldridge, 2010). Thứ hai, ở các nghiên cứu đơn quốc gia, luận án
bổ sung biến công cụ (instrumental variable) cho năng lực công nghệ, với
các kiểm định độ mạnh công cụ (first-stage F) và over-identification phù
hợp, nhằm tách hiệu ứng nhân quả khỏi nội sinh do lựa chọn và đồng thời
(P3, P5). Cần phân biệt nội sinh của năng lực, vốn được xử lý bằng biến
công cụ nói trên, với một mối đe dọa nhận dạng riêng là tự chọn lọc tham
gia xuất khẩu: theo logic Melitz (2003), doanh nghiệp năng suất cao tự
chọn vào xuất khẩu, nên hệ số FSTS có thể trộn lẫn hiệu ứng nhân quả của
quốc tế hóa với sự chọn lọc có sẵn. Luận án khoanh vùng đe dọa này theo
ba cách: (i) hiệu ứng cố định hai chiều cùng chuẩn hóa within
country--year hấp thụ chênh lệch mức năng suất cố định giữa doanh nghiệp
xuất khẩu và không xuất khẩu ở tầng nền kinh tế--năm; (ii) đối tượng suy
luận trọng tâm là hình dạng phi tuyến (độ cong và vị trí điểm uốn) chứ
không chỉ là mức hệ số, mà một câu chuyện chọn lọc mức đơn thuần khó tạo
ra một cực đại nội vùng có hệ thống; và (iii) nghiên cứu Singapore (P4)
bổ sung hiệu chỉnh chọn lọc kiểu Heckman với tỷ số Mills nghịch đảo như
một kiểm tra cục bộ. Tuy vậy, do hiệu chỉnh chọn lọc tham gia không được
thực hiện đồng bộ ở mọi nghiên cứu thành phần, luận án thừa nhận đây là
một ranh giới suy luận còn lại của thiết kế lát cắt ngang. Một kiểm tra
chọn lọc kiểu Heckman ở cấp gộp 50 nền (Mục 4.6.8) củng cố thêm: khi
thêm tỷ số Mills nghịch đảo từ probit tham gia xuất khẩu, độ cong chữ U
ngược và điểm uốn gần như không đổi và bản thân tỷ số Mills không có ý
nghĩa, cho thấy thiên lệch tự chọn lọc trong độ cong ước lượng là không
đáng kể trong miền dữ liệu. Thứ ba, độ ổn định của ước lượng trước thay
đổi do biến bỏ sót được đánh giá bằng giới hạn δ-stability theo Oster
(2019), một chuẩn mực ngày càng được áp dụng trong nghiên cứu thực
nghiệm khi không có thiết kế thực nghiệm thuần túy. Cuối cùng, luận án
tuân thủ nguyên tắc báo cáo trung thực: các kết quả từ mẫu nhỏ hoặc
thiếu công suất (ví dụ mẫu chỉ gồm doanh nghiệp xuất khẩu tại SIDS) được
trình bày như ``ranh giới suy luận'' (inferential bounds) chứ không phải
bằng chứng nhân quả khẳng định, và mọi suy luận về điều tiết thể chế
được neo đồng thời vào hai nguồn độc lập là kiểm định điều tiết của phân
tích tổng hợp và phổ điểm uốn của mô hình đa quốc gia (xem Mục 4.2.3,
Mục 4.6.2). Cách tiếp cận đa lớp này tuân theo các khuyến nghị về suy
luận từ dữ liệu quan sát quy mô lớn (Angrist \& Pischke, 2009). Một giới
hạn đo lường cần lưu ý là \textbf{tính bất biến đo lường (measurement
invariance)} của các cấu trúc TCI và DAI qua ba thế hệ bảng hỏi WBES: do
cấu trúc câu hỏi thay đổi giữa PICS3, Standardized và BREADY, một phần
khác biệt theo thời gian có thể phản ánh hiệu ứng lược đồ thay vì thay
đổi thực chất (ví dụ mức R\&D đo được ở một số nhóm). Luận án giảm thiểu
rủi ro này bằng hiệu ứng cố định năm/đợt, chuẩn hóa within
country--year, và kiểm tra độ nhạy với biến thể thước đo (Mục 3.5.1).
Luận án chưa thực hiện kiểm định bất biến đo lường đầy đủ (full
measurement invariance), nhưng một kiểm tra bất biến bộ phận (Mục 4.6.8)
cho thấy các mục lõi của TCI và DAI hiện diện cho khoảng \(99\%\) doanh
nghiệp ở cả ba thế hệ lược đồ và thay đổi trung bình theo thời gian phản
ánh khuếch tán số thực chất chứ không phải đứt gãy lược đồ; do đó các so
sánh xuyên-lược đồ trên các mục lõi này là có cơ sở, dù vẫn cần đọc thận
trọng.

\hypertarget{chiux1ebfn-lux1b0ux1ee3c-biux1ebfn-cuxf4ng-cux1ee5-vuxe0-kiux1ec3m-ux111ux1ecbnh-ux111ux1ed9-mux1ea1nh-cuxf4ng-cux1ee5}{%
\subsection{3.5.6.1 Chiến lược biến công cụ và kiểm định độ mạnh công
cụ}\label{chiux1ebfn-lux1b0ux1ee3c-biux1ebfn-cuxf4ng-cux1ee5-vuxe0-kiux1ec3m-ux111ux1ecbnh-ux111ux1ed9-mux1ea1nh-cuxf4ng-cux1ee5}}

Để tách hiệu ứng nhân quả của năng lực khỏi nội sinh do lựa chọn và đồng
thời, các nghiên cứu đơn quốc gia bổ sung ước lượng bình phương nhỏ nhất
hai bước (two-stage least squares, 2SLS). Công cụ được xây dựng từ tỷ lệ
áp dụng đồng cấp trong các ô ngành, vùng và đợt khảo sát theo nguyên tắc
loại bỏ chính doanh nghiệp khỏi trung bình ô (leave-one-out), nhằm bảo
đảm giá trị công cụ của một doanh nghiệp không chứa thông tin của chính
nó. Logic nhận dạng dựa trên hai điều kiện: tính liên quan, vì tỷ lệ áp
dụng đồng cấp dự báo mạnh quyết định áp dụng của doanh nghiệp riêng lẻ;
và loại trừ xấp xỉ, vì tỷ lệ áp dụng trung bình của các doanh nghiệp
khác trong cùng ô khó tác động đến năng suất của một doanh nghiệp qua
kênh nào khác ngoài chính quyết định áp dụng, sau khi đã kiểm soát ô
ngành và đợt.

Độ mạnh công cụ được đánh giá bằng thống kê F của hồi quy giai đoạn một,
với kết quả nằm trong khoảng 22 đến 35, vượt xa ngưỡng kinh nghiệm về
công cụ yếu. Một phát hiện phương pháp luận đáng chú ý là ước lượng 2SLS
cho hai chiều năng lực phân kỳ rõ rệt: hệ số công cụ hóa của năng lực
công nghệ dương mạnh (ước lượng 2SLS bằng 1,64, p \textless{} 0,001,
dưới một công cụ mạnh với F giai đoạn một bằng 22,1), trong khi hệ số
công cụ hóa của chấp nhận số ở mức Tầng 1 không khác biệt thống kê với
không, một kết quả độ vững củng cố cách diễn giải tính lỗi thời của biến
đại diện Tầng 1 trong môi trường số đã lan tỏa gần phổ quát.

\hypertarget{giux1edbi-hux1ea1n-ux1ed5n-ux111ux1ecbnh-hux1ec7-sux1ed1-oster-trux1b0ux1edbc-biux1ebfn-bux1ecf-suxf3t-khuxf4ng-quan-suxe1t}{%
\subsection{3.5.6.2 Giới hạn ổn định hệ số Oster trước biến bỏ sót không
quan
sát}\label{giux1edbi-hux1ea1n-ux1ed5n-ux111ux1ecbnh-hux1ec7-sux1ed1-oster-trux1b0ux1edbc-biux1ebfn-bux1ecf-suxf3t-khuxf4ng-quan-suxe1t}}

Vì dữ liệu chủ đạo là lát cắt ngang lặp lại chứ không phải thiết kế thực
nghiệm thuần túy, luận án đánh giá độ vững của hệ số trọng tâm trước
thiên lệch biến bỏ sót không quan sát bằng giới hạn ổn định delta của
Oster (2019). Phương pháp dựa trên trực giác rằng mức độ một hệ số dịch
chuyển khi thêm biến kiểm soát quan sát được mang thông tin về mức độ nó
có thể tiếp tục dịch chuyển nếu kiểm soát thêm các biến không quan sát
có cùng cấu trúc liên hệ. Tham số delta lượng hóa tỷ lệ tương đối giữa
lựa chọn trên biến không quan sát và lựa chọn trên biến quan sát; giá
trị delta bằng một biểu thị giả định bảo thủ rằng các biến không quan
sát quan trọng ngang với các biến quan sát.

Việc tính giới hạn đòi hỏi đặt cận trên cho mức độ giải thích tối đa của
một mô hình giả định đầy đủ biến. Luận án dùng quy ước cận trên này bằng
1,3 lần hệ số xác định của mô hình có kiểm soát, một mức được khuyến
nghị rộng rãi trong thực hành thực nghiệm. Kết luận phương pháp luận là
không hệ số trọng tâm nào đổi dấu hoặc sụp đổ về không dưới các độ lớn
khả tín của lựa chọn không quan sát, qua đó cung cấp một bằng chứng độc
lập rằng các mối liên hệ trọng tâm không phải là tạo phẩm của biến bỏ
sót có thể thấy trước.

\hypertarget{tuxednh-nhux1ea5t-quuxe1n-vux1ec1-dux1ea5u-qua-nhiux1ec1u-muxf4-huxecnh-nhux1b0-nguyuxean-tux1eafc-suy-luux1eadn-nux1ec1n-tux1ea3ng}{%
\subsection{3.5.6.3 Tính nhất quán về dấu qua nhiều mô hình như nguyên
tắc suy luận nền
tảng}\label{tuxednh-nhux1ea5t-quuxe1n-vux1ec1-dux1ea5u-qua-nhiux1ec1u-muxf4-huxecnh-nhux1b0-nguyuxean-tux1eafc-suy-luux1eadn-nux1ec1n-tux1ea3ng}}

Triết lý suy luận xuyên suốt khung độ vững là neo kết luận thực chất vào
tính nhất quán về dấu của ước lượng trọng tâm qua nhiều mô hình, chứ
không vào ý nghĩa của một hệ số đơn lẻ, vì các bảng độ vững kiểm tra các
quan ngại nhận diện khác nhau thay vì lặp lại cùng một giả thuyết.
Nguyên tắc này được minh họa rõ nhất ở nghiên cứu đảo nhỏ Thái Bình
Dương, nơi hệ số cường độ xuất khẩu giữ dấu âm nhất quán qua mọi mô
hình: âm và có ý nghĩa rất cao trong mô hình hiệu ứng cố định nền kinh
tế và năm với hệ số −1,339, vẫn âm và có ý nghĩa khi chỉ giữ hiệu ứng cố
định năm hoặc khi bỏ toàn bộ biến kiểm soát, và chỉ mất ý nghĩa ở mẫu
con chỉ gồm doanh nghiệp xuất khẩu vốn quá nhỏ. Suy luận về gánh nặng
quốc tế hóa bắt buộc do đó được neo chủ yếu vào các mô hình hiệu ứng cố
định chứ không vào mẫu con thiếu công suất, đúng nguyên tắc báo cáo
trung thực đã nêu ở Mục 3.5.6.

Nguyên tắc tương tự áp dụng cho các kết quả điều tiết không có ý nghĩa
trong khung đa quốc gia. Năng lực công nghệ vào mô hình với hiệu ứng mức
dương và rất có ý nghĩa, cho thấy doanh nghiệp năng lực cao hơn năng
suất hơn ở mọi mức quốc tế hóa, nhưng các số hạng tương tác với cường độ
xuất khẩu đều không có ý nghĩa và điểm uốn gần như không đổi, nghĩa là
năng lực nâng nền mà không bẻ cong đường. Kết quả không có ý nghĩa của
các số hạng tương tác vì vậy không bị diễn giải là thất bại đo lường mà
là bằng chứng tích cực cho cách phân vai năng lực công nghệ là bộ tăng
mức chứ không phải bộ điều tiết độ cong, một kết luận được củng cố thêm
bởi sự nhất quán của hiệu ứng mức dương của năng lực công nghệ ở các
nghiên cứu đơn quốc gia. Việc đối chiếu chéo giữa hiệu ứng mức ổn định
và hiệu ứng điều tiết không có ý nghĩa, lặp lại qua nhiều nghiên cứu
thành phần và nhiều chế độ thể chế, là chính cấu trúc bằng chứng cho
phép luận án phân biệt một quy luật cấu trúc bền vững khỏi một hiệu ứng
ngẫu nhiên của mẫu.

\hypertarget{tuxf3m-tux1eaft-phux1b0ux1a1ng-phuxe1p-kux1ebf-thux1eeba-vuxe0-ux111uxf3ng-guxf3p-mux1edbi}{%
\section{3.6 Tóm tắt phương pháp: Kế thừa và đóng góp
mới}\label{tuxf3m-tux1eaft-phux1b0ux1a1ng-phuxe1p-kux1ebf-thux1eeba-vuxe0-ux111uxf3ng-guxf3p-mux1edbi}}

\begin{longtable}[]{@{}
  >{\raggedright\arraybackslash}p{(\columnwidth - 4\tabcolsep) * \real{0.3333}}
  >{\raggedright\arraybackslash}p{(\columnwidth - 4\tabcolsep) * \real{0.3333}}
  >{\raggedright\arraybackslash}p{(\columnwidth - 4\tabcolsep) * \real{0.3333}}@{}}
\toprule\noalign{}
\begin{minipage}[b]{\linewidth}\raggedright
Thành phần
\end{minipage} & \begin{minipage}[b]{\linewidth}\raggedright
Cơ sở kế thừa
\end{minipage} & \begin{minipage}[b]{\linewidth}\raggedright
Đóng góp mới
\end{minipage} \\
\midrule\noalign{}
\endhead
\bottomrule\noalign{}
\endlastfoot
Thiết kế hỗn hợp & Borenstein et al.~(2009); Hunter \& Schmidt (2004) &
Áp dụng cho \textbf{50 nền kinh tế châu Á + Thái Bình Dương} trong một
luận án \\
phân tích tổng hợp 1977--2026 & Bausch \& Krist (2007); Kirca et
al.~(2012); Marano et al.~(2016) & Mở rộng phạm vi, bổ sung digital và 6
chế độ con ICRV moderators \\
Mẫu thực nghiệm gộp & nghiên cứu thành phần P1, 17 nền kinh tế châu Á
mới nổi & Mở rộng từ 17 lên \textbf{50 nền kinh tế}, 103 cặp nền kinh tế
và năm (gồm Nhật Bản 2025) \\
Biến phụ thuộc năng suất lao động & Bloom et al.~(2012); Hsieh \& Klenow
(2009); nghiên cứu thành phần P1 & Lập luận rõ về hiệu quả hoạt động
kinh doanh của doanh nghiệp đa chiều trong bối cảnh WBES \\
FSTS + FSTS² & Lu \& Beamish (2004); Hitt et al.~(1997); nghiên cứu
thành phần P2 và P5 & Đưa vào cùng khung điều tiết ba chiều \\
TCI vs DAI & Bhandari et al.~(2023); Verhoef et al.~(2021) & Không chồng
lấn, độ thuần cấu trúc protocol cho WBES; lá chắn số mở rộng (nghiên cứu
thành phần P1) \\
Lind--Mehlum U-test & Lind \& Mehlum (2010); Haans et al.~(2016) & Áp
dụng cho từng country sample và multi-country \\
Điều tiết ba chiều & Aiken \& West (1991); Dawson (2014) & số × thể chế
× nhà quản trị, mới cho châu Á và Thái Bình Dương \\
dị biệt theo thời gian & Wu et al.~(2022); Xiao et al.~(2013); nghiên
cứu thành phần P2 & kiểm định ổn định Trung Quốc 2012--2024 (điểm uốn
\textasciitilde47,8\% \(\mathrm{FSTS}\)) \\
sai số chuẩn vững & Long \& Ervin (2000); White (1980) & Chuẩn \\
lệch lạc công bố & Egger et al.~(1997); Begg \& Mazumdar (1994) &
Chuẩn \\
\end{longtable}

\hypertarget{ux111ux1ea1o-ux111ux1ee9c-nghiuxean-cux1ee9u-vuxe0-khux1ea3-nux103ng-tuxe1i-lux1eadp}{%
\section{3.7 Đạo đức nghiên cứu và khả năng tái
lập}\label{ux111ux1ea1o-ux111ux1ee9c-nghiuxean-cux1ee9u-vuxe0-khux1ea3-nux103ng-tuxe1i-lux1eadp}}

Luận án sử dụng dữ liệu thứ cấp từ WBES, bộ dữ liệu công khai được World
Bank cung cấp miễn phí cho mục đích nghiên cứu (World Bank Enterprise
Surveys, 2025). Tất cả các trích dẫn theo APA 7th để ghi nhận đầy đủ
công sức của các tác giả trước, tránh đạo văn. Code và dữ liệu phân tích
được lưu trữ trên các nhánh chuyên biệt của kho mã để bảo đảm khả năng
tái lập kết quả.

Về bảo mật và quyền riêng tư, dữ liệu WBES đã được World Bank ẩn danh ở
nguồn: bộ dữ liệu vi mô không chứa tên doanh nghiệp hay thông tin nhận
dạng cá nhân của đáp viên, và mọi định danh chỉ là mã số nội bộ phục vụ
ghép nối. Luận án không thu thập dữ liệu sơ cấp từ con người nên không
phát sinh nghĩa vụ xin chấp thuận đạo đức từ hội đồng đối với đối tượng
nghiên cứu là người; tuy vậy luận án vẫn tuân thủ điều khoản sử dụng của
World Bank, không thực hiện bất kỳ thao tác nào nhằm tái nhận dạng doanh
nghiệp, và chỉ báo cáo kết quả ở cấp tổng hợp hoặc cấp hệ số. Việc xử lý
mã không phản hồi của bộ công cụ, gồm mã không biết, mã từ chối và mã
không áp dụng, thành giá trị khuyết được thực hiện minh bạch theo quy
ước nêu ở Phụ lục A nhằm tránh diễn giải sai các mã kỹ thuật thành giá
trị thực.

Về khả năng tái lập, luận án áp dụng nguyên tắc minh bạch kiểu PRISMA
cho dữ liệu thứ cấp: toàn bộ quy trình hợp nhất sáu bước, từ lọc loại
khảo sát đến tạo mẫu phân tích bằng xóa theo danh sách, được ghi rõ tiêu
chí và số quan sát còn lại ở từng bước để cho phép tái dựng và kiểm
chứng, cùng dòng dữ liệu kiểu PRISMA và mã tái lập trình bày đầy đủ ở
Phụ lục A. Các con số trọng tâm của khung thực nghiệm được khóa và có
thể tái lập đúng, gồm mô hình bậc hai gốc với cỡ mẫu hồi quy 81.022 và
điểm uốn 51,5\%, cùng mô hình đầy đủ thêm năng lực công nghệ và chấp
nhận số với cỡ mẫu 79.080 và điểm uốn 43,6\%. Tập lệnh xây dựng dữ liệu
gộp tái lập đúng quy trình sáu bước S1 đến S6 trên kho mã hiện hành, bảo
đảm rằng mỗi cập nhật bổ sung nền kinh tế hay đợt khảo sát mới, chẳng
hạn Nhật Bản 2025, đều đi qua cùng một chuỗi quyết định lọc và hài hòa
hóa có thể kiểm chứng.

\end{document}
